\documentclass[11pt,a4paper]{article}
\pdfmapfile{+lm.map}
\pdfmapfile{+cm.map}
\pdfmapfile{+cmextra.map}
\pdfmapfile{+symbols.map}
\pdfmapfile{+euler.map}
\usepackage[T1]{fontenc}
\usepackage[utf8]{inputenc}
\usepackage{lmodern}
\usepackage[margin=27mm]{geometry}
\usepackage{amsmath,amssymb,amsthm,mathtools}
\usepackage{booktabs,longtable,array,enumitem}
\usepackage{xcolor}
\usepackage{microtype}
\usepackage{xurl}
\usepackage[colorlinks=true,linkcolor=blue!45!black,citecolor=blue!45!black,urlcolor=blue!45!black]{hyperref}
\newtheorem{theorem}{Theorem}[section]
\newtheorem{proposition}[theorem]{Proposition}
\newtheorem{lemma}[theorem]{Lemma}
\newtheorem{corollary}[theorem]{Corollary}
\newtheorem{sourcefact}[theorem]{Cited source statement}
\theoremstyle{definition}
\newtheorem{definition}[theorem]{Definition}
\newtheorem{remark}[theorem]{Remark}
\setlength{\emergencystretch}{2em}
\setlength{\parskip}{3pt}
\title{K17: Step 17, the $K$-groups of $\mathbb Z$,\\
Soul\'e characters and the pro-$\ell$ kernels of the tripod}
\author{Research run K17---checked deductions and explicit open gates}
\date{21 November 2026\\Timed checkpoint; research in progress}
\begin{document}
\maketitle
\begin{abstract}
The exact quantitative Selmer theorem behind step 17 is recovered from
the pinned source. Its rational-field specialization, together with
etale $K$-theory and arithmetic duality, bounds the orders of
$K_{4k}(\mathbb Z)$; a reconstruction at fixed levels yields exponent
seven. We separate the exact Vandiver equivalence, full-domain Soul\'e
indices, and the stronger integral generation claims. Complete
deductions settle the pro-$2$ kernel-field equality with unrestricted
real places and every odd prime below $1000$. A general sufficient
criterion uses Vandiver and odd cup-product surjectivity, with no
simple-root hypothesis. Integral odd generation fails at $37$ in
weight $32$ and at $691$ in weight $12$. At $2$ it holds through
weight $18$ by an exact compact-group computation; the all-weight
assertion and the remaining irregular-prime cases retain explicit gates. Cote 149, scans
79--85, is read through exact limiting Kummer sequences, with a concrete
completed-unit class outside the geometric tangent image.

This checkpoint is not the final package. Every completed new proof
has a whole-proof reading recorded in the check ledger. Literature
theorems, exact arithmetic computations and source inspections have
distinct check types. No new Lean build, formal certificate or external
peer review is claimed. Unresolved assertions appear separately.
\end{abstract}
\tableofcontents
\clearpage
\part{Checked deductions and rechecked source statements}

\section{Recovering the arithmetic statement behind step 17}
\label{sec:step17}

\paragraph{Source contract and provenance.}
Alp\"oge's \emph{Remark}, p.~4, step~17 and footnote~6, names
$H^1(\mathcal O_E[1/\ell],\mathbb Q_\ell/\mathbb Z_\ell(-e))$
and a weakening of Vandiver, but supplies neither quantifiers nor a theorem
number for that weakening. The attached map identifies candidate modules only.
We recovered those modules from the public mirror
\url{https://github.com/CocoSgt/riemann-hypothesis-lean}, at commit
\texttt{7cbf7a0c32fca1badf05931879a37c3f8a2b03b4}.
The retrieved \path{src/Main.lean} is byte-identical to the supplied file;
its SHA-256 is recorded in the evidence manifest. Git blob identifiers were
checked for the manifest-based downloads. These are source inspections, not
a new Lean build or a claim to have authenticated the entire mirror. Under
the prompt's stipulated theorem contract, we use the recovered unconditional
Alp\"oge theorem below. The arithmetic translation and consequences are the
new deductions checked here.

In this section a source path beginning with \texttt{AV/} abbreviates
\path{src/Development/AutomorphicArithmetic/AutomorphicVarieties/}.
The source facts have precise declaration locators; no theorem is inferred
merely from a file name or from a proposition's definition.

\subsection{The local conditions and the fields}

Put $D_\ell(e)=\mathbb Q_\ell/\mathbb Z_\ell(-e)$ and, for a number field $E$,
\[
 \mathcal S_\ell(E,e)=
 \ker\left(H^1_{\mathrm{et}}(\operatorname{Spec}\mathcal O_E[1/\ell],D_\ell(e))
       \longrightarrow\prod_{v\mid\ell}H^1(E_v,D_\ell(e))\right).
 \tag{S1}
\]
Thus the selected classes are unramified outside $\ell$ and locally trivial
at every place above $\ell$. They are a subgroup of the displayed $H^1$,
not that entire group. The source definitions are
\path{src/Formalism/GaloisCohomology/Statement.lean}, declarations
\texttt{LData}, \texttt{qzTw}, \texttt{Sha}, and \texttt{shaLen}.

For a finite extension coefficient ring $O\subset\overline{\mathbb Q}_\ell$
which is finite over $\mathbb Z_\ell$ and is a DVR, write
\[
 D_O=O\otimes_{\mathbb Z_\ell}(\mathbb Q_\ell/\mathbb Z_\ell),
 \qquad e(O)=\operatorname{length}_O(O/\ell O).
 \tag{S2}
\]
For a Dirichlet character $\psi$ modulo $N$ with values in $O^\times$, define
$\mathcal V_{\ell,N}(e,\psi)$ to be the subgroup of
$H^1(G_{\mathbb Q},D_O(\psi^{-1}\chi_\ell^{-e}))$
whose restriction is zero on every finite inertia group and on every
decomposition group at a prime dividing $\ell N$. Here $\chi_\ell$ is the
cyclotomic character. Set
\[
 v_{\ell,N}(e,\psi)
 =\operatorname{length}_O\mathcal V_{\ell,N}(e,\psi)/e(O).
 \tag{S3}
\]
In (S3) the quotient of lengths is used only when the module has finite
length. Literally, the source's \texttt{Vn} applies
\texttt{ENat.toNat} to its extended-natural lengths (which sends
infinity to zero); every estimate used below simultaneously proves
finiteness, so (S3) then gives exactly its value. These are
\texttt{Coef}, \texttt{Vlen}, and \texttt{Vn} in
\path{AV/HeightAndVanishing/DirichletPrimeToLPartMulLPowerPart.lean}, together
with \texttt{QO}, \texttt{eO}, \texttt{shaFam}, and \texttt{ShaU} in
\path{AV/GaloisLibrary/SelmerTransport/ConjStableShaFam.lean}.
The scalar character is defined in
\path{AV/CharacterLevelBounds/ShaFamMono.lean}, \texttt{rho}.

\begin{proposition}[K17-FIELD-DESC: the field interpretation]
Let $e>0$, let $\ell$ be odd, and suppose $E/\mathbb Q$ is a cyclic
extension of degree prime to $\ell$, cut out by $\psi$. Suppose $O$ contains
the values of $\psi$. Restriction gives an isomorphism
\[
 \mathcal V_{\ell,N}(e,\psi)
 \simeq
 \left(O\otimes_{\mathbb Z_\ell}\mathcal S_\ell(E,e)\right)^{\psi},
 \tag{S4}
\]
where the right side is the $\psi$-eigenspace for the natural action of
$\operatorname{Gal}(E/\mathbb Q)$ on cohomology, including its action on
the Tate coefficients. The same description applies to character components
inside a larger abelian extension of degree prime to $\ell$.
\end{proposition}
\begin{proof}
Write $\Delta=\operatorname{Gal}(E/\mathbb Q)$. All positive-degree
cohomology of $\Delta$ with $\ell$-primary coefficients vanishes: it is
annihilated by $|\Delta|$, which acts invertibly on those coefficients.
Inflation--restriction consequently identifies the global $H^1$ on the
left with the indicated eigenspace of $H^1(G_E,D_O(-e))$. The same argument
applies to decomposition and inertia groups, since their quotient groups
have order dividing $|\Delta|$. Hence restriction detects the vanishing of
each local restriction.

An unramified local class at a place $w\nmid\ell$ of $E$ lies in
\[
 H^1_{\mathrm{unr}}(E_w,D_O(-e))
 =D_O/(\operatorname{N}w^{-e}-1)D_O=0.
\]
The last equality holds because $\operatorname{N}w^{-e}-1$ is a nonzero
element of $\mathbb Q_\ell$ and multiplication by any nonzero scalar from
$\mathbb Z_\ell$ is surjective on the divisible module $D_O$.
Thus all extra decomposition conditions away from $\ell$, including
those dividing $N$, are automatic after the inertia conditions. The local
conditions at $\ell$ are exactly (S1). Restriction therefore gives (S4).
Finite flatness of $O$ over $\mathbb Z_\ell$ commutes with the cochain
description and the eigenspace idempotent, so no scalar-extension defect
remains. This proves the assertion.
\end{proof}

For CM fields the source's ``plus'' notation is twist-dependent. If $c_E$
is complex conjugation, then
\[
 \mathcal S_\ell(E,e)^+
 =\{x\in\mathcal S_\ell(E,e):c_E x=(-1)^e x\}.
 \tag{S5}
\]
It is not always the $+1$ eigenspace. See
\path{AV/LengthEqAddEigen.lean}, \texttt{ShaPlus}, \texttt{ShaMinus},
\texttt{shaLen\_eq\_add}, and
\path{AV/PartsSign.lean}, \texttt{mk\_mem\_ShaPlus\_iff}.
The direct sum of these two eigenspaces uses that $2$ is invertible;
the source treats $\ell=2$ separately.

The fields used in the larger twisting argument are not arbitrary number
fields. The source's \texttt{Host} fixes a Dirichlet character $\chi$ of
modulus $m$, an integer $a\ge2$, and primes $q\equiv3\pmod8$,
$d\equiv7\pmod{16}$ with $(qd,m)=1$. They determine
\[
 E=F_\chi\,\mathbb Q(\mu_{2^{a+2}})^+
      \mathbb Q(\sqrt{-q})\,\mathbb Q(\sqrt{-d}),
 \qquad N=\operatorname{lcm}(2^{a+2},m)qd,
 \tag{S6}
\]
where $F_\chi$ is the fixed field of $\ker\chi$. This is the fixed-field
description encoded by \texttt{IsCompField} and
\texttt{isCompField\_ofChar} in \path{AV/StmtFields.lean}, and by
\texttt{Host.H} in
\path{AV/HeightAndVanishing/DirichletPrimeToLPartMulLPowerPart.lean}.
Indeed the four congruence conditions in that definition fix respectively
the real $2$-power cyclotomic field, $F_\chi$, and the two quadratic fields;
the intersection of the corresponding subgroups fixes their compositum.
The selected character components have $\psi(-1)=(-1)^e$.

For completeness, the exact finite-parameter organization is in
\path{AV/HeightAndVanishing/LayerLe.lean}, where \texttt{PointHV}
is the proposition defined to assert the following statement:
for each $\chi$ and function $j:\mathbb N\to\mathbb N$, there are
$C,r_1,r_2\in\mathbb N$, with $r_1,r_2$ prime, such that, for every host
with character $\chi$, both auxiliary primes smaller than $q,d$, every
$t\ge2$ satisfying \texttt{Host.Big}, and every $B\ge0$ satisfying
the specified \texttt{LayerLe} bounds at the finite set \texttt{Bchi},
the weighted sum of the lengths in (S5), over all odd primes and
$1\le e<t$, is bounded by $(Nt)^C+CB$.
Here \texttt{Host.Big} means $2t+1<q,d$ and each of
$m,\operatorname{ord}\chi,p_K(\chi),h_K(\chi),p_1(\chi)$ is less than
$q,d$. The finite exceptional set is
\[
 B_\chi=\{3\}\cup
 \{\text{odd prime divisors of }
   (\operatorname{ord}\chi)m p_K(\chi)h_K(\chi)p_1(\chi)\}.
\]
The three auxiliary quantities are the least primes/class number specified
in \path{AV/StmtFields.lean}, \texttt{pK}, \texttt{hStar}, \texttt{pOne}.
We record this definition to identify the fields and classes in the
note. The closed source proof composition in the following host-family
subsection supplies its theorem inputs; the stated layer bounds remain
explicit premises. The deductions about $\mathbb Z$ below use the unconditional,
quantified theorem next, and do not assume an unverified numerical value
for any of these constants or layer bounds.

\subsection{The unconditional character theorem}

The exact admissibility condition \texttt{GD} is in
\path{AV/HeightAndVanishing/Counts.lean}. In ordinary notation it says:
$e\ge1$, $N$ is even, $\ell$ is odd, the order of $\psi$ is prime to
$\ell$, the factor $\psi_\ell$ supported at the $\ell$-part of the
conductor satisfies $\psi_\ell^2=1$, and
$\psi(-1)=(-1)^e$. There is an auxiliary imaginary quadratic field $K$
such that
\[
 \ell\nmid |D_K|h_K,\quad (|D_K|,N)=1,\quad
 K\not\subset\mathbb Q(\mu_\ell)F_{\psi/\psi_\ell},\quad
 (\ell-1)\nmid2e. \tag{S7}
\]
The third condition is exactly the source's subgroup noncontainment
\texttt{eta}; it has not been silently deleted. The source definition
of \texttt{chiP} agrees with $\psi/\psi_\ell$ under its prime-to-$\ell$
order hypothesis. In particular these conditions force $\ell\ge5$.

\begin{sourcefact}[Alp\"oge: K17-S17, the recovered quantitative theorem]
For every imaginary quadratic $K$ there exist $c_5,A\in\mathbb N$ with
the following property. Fix $e,N$, a subgroup
$H\subset(\mathbb Z/N\mathbb Z)^\times$, and distinct primes $p,p'$
not dividing $30N|D_K|$. For each prime $\ell$ in a finite set $L$ choose
a permitted coefficient ring $O_\ell$. If at least one character with
kernel $H$ satisfies (S7) and the preceding admissibility conditions at
each $\ell\in L$, then all the associated lengths are finite and
\[
 \sum_{\ell\in L}\log\ell
   \sum_{\ker\psi=H}v_{\ell,N}(e,\psi)
 \le 2c_5\varphi([(\mathbb Z/N\mathbb Z)^\times:H])
          (N(e+1))^A(\log p+\log p'). \tag{S8}
\]
The character sum is over the characters valued in $O_\ell$; existence
of one such faithful character provides its whole orbit.

Exact locators: \path{AV/UnitaryThreeGroups/KatoUse.lean},
\texttt{CorGK}; and \path{AV/SmallClassNumbers/CorGKOf.lean},
\texttt{UpperTheory.V.SmallClassNumbers.corGK}, an unconditional theorem.
\end{sourcefact}

\paragraph{Proof reconstruction and the cited inputs.}
The arithmetic input \texttt{PropGK} supplies, for an auxiliary prime $p$,
a nonzero integer $n_p$ such that
\[
 \log|n_p|\le c_5\varphi([G:H])(N(e+1))^A\log p,
 \qquad
 \sum_{\ker\psi=H}v_{\ell,N}(e,\psi)\le2v_\ell(n_p)
 \quad(\ell\ne p). \tag{S9}
\]
Its source proof is \texttt{propGK} in
\path{AV/UnitaryThreeGroups/OrbEquiv.lean}. It combines
\texttt{OneChar} (the Selmer-length estimate by rank-three/rank-two
automorphic congruences) with \texttt{CongruenceNumberRankThree}
(a nonzero norm of Frobenius-trace differences). The latter provides the
integer and its archimedean size, while the former bounds its local
valuation. The faithful-character orbit has size $\varphi([G:H])$.
Their unconditional inputs are assembled in
\texttt{SmallClassNumbers.corGK\_of}, \texttt{litN}, and \texttt{corGK}
of the file cited above: local and global Galois cohomology, the Kato
input, adjoint vanishing, automorphy, and congruence-number estimates.
These are cited Alp\"oge theorems under the source contract; they are
not independent new automorphy proofs in this paper.

Here is the complete elementary passage from (S9) to (S8).
Use (S9) twice, with $p$ and $p'$. For each $\ell$ at least one applies,
so the nonnegative length sum is at most
$2v_\ell(n_p)+2v_\ell(n_{p'})$. For a nonzero integer $n$,
\[
 \sum_{\ell\in L}v_\ell(n)\log\ell\le\log|n|
\]
because the product of the indicated prime powers divides $|n|$.
Summing and using the two archimedean estimates proves (S8), including
the auxiliary-prime cases. This is precisely
\path{AV/UnitaryThreeGroups/CorGSumLog.lean}, \texttt{corGK}.
The complete source proofs of this passage and of the orbit/norm
translation were read in this run.

\subsection{The rational-field specialization}

\begin{theorem}[K17-S17-Q]
There exist absolute $C,A\in\mathbb N$ such that for every positive even
integer $e$ and every finite set of primes $L$ satisfying
$\ell\ge5$ and $(\ell-1)\nmid2e$,
\[
 \mathcal S_\ell(\mathbb Q,e)\text{ is finite},\qquad
 \sum_{\ell\in L}\log\ell\,
   \operatorname{length}_{\mathbb Z_\ell}\mathcal S_\ell(\mathbb Q,e)
 \le C(e+1)^A. \tag{S10}
\]
In particular the group is zero whenever $\log\ell>C(e+1)^A$
and the displayed prime conditions hold.
\end{theorem}
\begin{proof}
In (S8) take $K=\mathbb Q(\sqrt{-3})$, $N=2$, $H=G$, the trivial
character, and auxiliary primes $p=7,p'=11$. The discriminant is $-3$
and the class number is one. The character conditions are immediate;
evenness of $e$ gives the parity condition. For $\ell\ge5$ the field
$K$ is not contained in $\mathbb Q(\mu_\ell)$, since the latter is
unramified at $3$ whereas $K$ is ramified there. This checks (S7).
All these checks are also made by \texttt{gd} in
\path{AV/SmallClassNumbers/ClassNumberThree.lean}.

Choose the bottom coefficient ring. It is isomorphic to
$\mathbb Z_\ell$ by \texttt{exists\_coef} in
\path{AV/SmallClassNumbers/LenEqZeroIff.lean}. Then $D_O=D_\ell(0)$
and $e(O)=1$, since $O/\ell O\simeq\mathbb F_\ell$ is a simple
$O$-module. The extra local condition at $2$ is automatic by the
unramified local calculation in the proof of K17-FIELD-DESC.
Thus the length in (S8) is the length in (S10).

The right side of (S8) is
$2c_5 2^A(e+1)^A(\log7+\log11)$, bounded by
$36c_5 2^A(e+1)^A$. Take $C=36c_5 2^A$.
This is the proof \texttt{bound\_of} in
\path{ClassNumberThree.lean}; its unconditional input and conclusion
are \texttt{corGK} and \texttt{bound} in
\path{SmallClassNumbers/CorGKOf.lean}.
For the last assertion use $L=\{\ell\}$. A nonzero finite
$\mathbb Z_\ell$-module has length at least one, contradicting
the strict inequality on $\log\ell$. The source independently states
this cocycle formulation as \texttt{Large} and proves
\texttt{large\_of} in \path{LenEqZeroIff.lean}.
\end{proof}

\paragraph{What is and is not numerical here.}
The source has also a proposition \texttt{CorG} with a displayed
$e^7N^{18}|D_K|^5$ shape. A definition of that proposition is not a
proof of it. We have used the established \texttt{CorGK}, with its
existential exponent, throughout. Neither ``100ish'' in footnote~6
nor an unspecified constant is a certified degree range. The next
section gives every vanishing and cyclicity conclusion that actually
follows from (S10), and identifies the remaining numerical gate.

\subsection{A closed theorem for the Case~I host family}
\label{sec:host-case-one}

\paragraph{K17-HOST-CASE-I: a theorem, beyond the field definitions.}
There is a closed source proof of the host-family assertion
\texttt{CountI}.  In full namespace notation its proof term is the
composition
\[
 \texttt{UpperTheory.V.U3.countI}\;
 \texttt{UpperTheory.V.SmallClassNumbers.corGK}.
\]
The first theorem has type
\(\texttt{U3.CorGK}\longrightarrow\texttt{HV.CountI}\);
the second is the closed theorem already used for the rational-field
specialization.  This composition uses the established, field-dependent
polynomial theorem \texttt{CorGK}.  It does not require the stronger
uniform discriminant and level exponents appearing in the separately
defined proposition \texttt{HV.CorG}.

Here is the precise mathematical reading.  Fix a character \(\chi\) in
the source's \texttt{PChar} domain, of modulus \(m\), and put
\(o=\operatorname{ord}(\chi)\).
Let \(p_K\) be the least prime at least seven congruent to three modulo
four and coprime to \(m\); set
\[
 K_\chi=\mathbb Q(\sqrt{-p_K}),\qquad
 h_\chi=|\operatorname{Cl}(K_\chi)|.
\]
Let \(p_1\) be the least prime congruent to one modulo the odd part of
\(o\), and coprime to \(2mp_K\).  The finite set used by the source is
\[
 \mathcal B_\chi=\{3\}\ \cup\
 \{\ell:\ell\text{ is an odd prime dividing }omp_Kh_\chi p_1\}.
\]
There exist natural numbers \(C\) and distinct primes \(r_1,r_2\),
depending only on \(\chi\), with the following property.

Choose a host \(h=(\chi,a,q,d)\), where \(a\ge2\), \(q,d\) are primes,
\(q\equiv3\pmod8\), \(d\equiv7\pmod{16}\), and
\(\gcd(qd,m)=1\).  These are exactly the source's
\texttt{PrimePair} conditions.  Its level and associated field are
\[
 N=\operatorname{lcm}(2^{a+2},m)qd,
 \qquad
 E_h=F_\chi\,\mathbb Q(\mu_{2^{a+2}})^+
          \mathbb Q(\sqrt{-q})\,\mathbb Q(\sqrt{-d}).
\]
Require
\[
 r_i<q,\quad r_i<d\quad(i=1,2),
\]
and, for \(t\in\mathbb N\), the explicit host-size conditions
\[
 2t+1<q,d,
 \qquad y<q,d\quad
 (y\in\{m,o,p_K,h_\chi,p_1\}).
\]
These are precisely \(h.\texttt{Big}(t)\).  Let \(e\in\mathbb N\) with \(0<e<t\), and let
\(L\) be any finite set of odd primes such that
\[
 \ell\notin\mathcal B_\chi,
 \qquad (\ell-1)\nmid 2e
 \qquad(\ell\in L).
\]
This pair of conditions is exactly \(h.\texttt{CaseI}(\ell,e)\).

For each \(\ell\), take any coefficient ring \(\mathcal O_\ell\) of the
source, and let \(X_h(\mathcal O_\ell,e)\) be the characters modulo
\(N\) that factor through \(\operatorname{Gal}(E_h/\mathbb Q)\) and
satisfy \(\psi(-1)=(-1)^e\).  Use \(V_{\ell,N}(e,\psi)\) for the
cohomology group with the everywhere-inertia and selected decomposition
conditions defined in the source reconstruction above.  Then all these
groups, for \(\ell\in L\) and \(\psi\in X_h(\mathcal O_\ell,e)\),
have finite length, and
\begin{equation}\label{eq:host-case-one-bound}
 \sum_{\ell\in L}\log\ell\!
 \sum_{\psi\in X_h(\mathcal O_\ell,e)}
 \frac{\operatorname{length}_{\mathcal O_\ell}
                   V_{\ell,N}(e,\psi)}
      {\operatorname{length}_{\mathcal O_\ell}
                   (\mathcal O_\ell/\ell\mathcal O_\ell)}
 \le (Nt)^C.
\end{equation}
The constants precede the host, the twist, the coefficient choices and
the finite prime set.  Both even and odd \(e\) are permitted here;
the character parity changes with \(e\).  No assertion about the whole
ambient \(H^1\) is substituted for these local kernels.

\paragraph{Proof and source check.}
Apply \texttt{CorGK} to the fixed imaginary quadratic field
\(K_\chi\), obtaining constants \(c_5,A\).  Choose
\(r_1\ge30mp_K+1\) and \(r_2\ge r_1+1\), both prime.
The prime divisors of the host level are among those dividing \(2mqd\).
The displayed inequalities therefore ensure
\(r_i\nmid30N|\operatorname{disc}(K_\chi)|\).
The source theorem \texttt{gdOfHost}, with its complete proof checked,
shows that every character in \(X_h(\mathcal O_\ell,e)\) satisfies the
\texttt{GD} hypotheses for this fixed \(K_\chi\).  Its field exclusion
is proved by ramification at \(p_K\): this prime ramifies in
\(K_\chi\), while the cyclotomic and character fields occurring in the
opposing compositum are unramified there.  The class-number and character
order exclusions are exactly supplied by \(\ell\notin\mathcal B_\chi\).

Group the characters by the kernel \(H\) of their character on
\((\mathbb Z/N\mathbb Z)^\times\).  Each quotient by such a kernel is
cyclic.  Applying \texttt{CorGK} to each resulting character orbit gives
finiteness and a bound with orbit factor \(\varphi([G:H])\), where
\(G=(\mathbb Z/N\mathbb Z)^\times\).  The source's finite character
count proves
\[
 \sum_H\varphi([G:H])\le |G|=\varphi(N)\le N.
\]
Indeed the characters with a fixed kernel are the faithful characters
of that cyclic quotient, numbering \(\varphi([G:H])\), and different
kernels give disjoint sets of characters.  Over a characteristic-zero
algebraically closed coefficient field there are exactly \(|G|\)
characters of the finite abelian group \(G\).  This also proves the
inequality without assuming the individual rings \(\mathcal O_\ell\)
already contain every root of unity.

Consequently the left side of \eqref{eq:host-case-one-bound} is at most
\[
 2c_5(\log r_1+\log r_2)\,N\,[N(e+1)]^A.
\]
Choose \(k\in\mathbb N\) with
\(2c_5(\log r_1+\log r_2)\le16^k\).
Since \(N\ge16\), \(e+1\le t\), and \(t\ge1\), this is at most
\((Nt)^{k+1+A}\).  Set \(C=k+1+A\).
These are the choices and inequalities in the checked source proof.

The exact implementation locators, relative to
\path{src/Development/AutomorphicArithmetic/AutomorphicVarieties/}, are
\path{UnitaryThreeGroups/CountI.lean}, theorem
\texttt{UpperTheory.V.U3.countI};
\path{SmallClassNumbers/CorGKOf.lean}, theorem
\texttt{UpperTheory.V.SmallClassNumbers.corGK};
\path{HeightAndVanishing/GdOfHost.lean}, theorem \texttt{gdOfHost};
\path{HeightAndVanishing/Counts.lean}, definition \texttt{CountI};
and \path{HeightAndVanishing/DirichletPrimeToLPartMulLPowerPart.lean},
definitions \texttt{Host}, \texttt{Big}, \texttt{X}, and \texttt{CaseI}.
The auxiliary primes and bad set are defined in \path{StmtFields.lean}.
The complete \texttt{U3.countI} proof, including its quantifier order,
auxiliary-prime exclusions, finite orbit partition and final absorption
inequality, was independently read.  No fresh Lean compilation is
claimed.

\paragraph{K17-HOST-POINT: the full point assertion also closes.}
Further retrieval supplies the remaining inputs to
\texttt{pointHV\_of\_counts}.  Consequently \texttt{PointHV} has a
closed proof by composition of the recovered source theorems.  Its
arithmetic layer-bound premise remains part of the theorem; closing
the source proof does not set that bound to zero.

For precision, fix \(\chi\) and any function
\(j:\mathbb N\longrightarrow\mathbb N\).  There are
\(C\in\mathbb N\) and primes \(r_1,r_2\), depending only on
\(\chi,j\), such that the following holds.  Take any host \(h\) as
above, with \(r_i<q,d\), any \(t\ge2\) satisfying
\(h.\texttt{Big}(t)\), and any real number \(B\ge0\).
For each \(\ell\in\mathcal B_\chi\), assume the following layer
condition: for every cyclotomic \(\ell\)-power extension
\[
 F/\mathbb Q(\mu_{\ell m}),\qquad
 [F:\mathbb Q(\mu_{\ell m})]=\ell^{j(\ell)},
\]
all the groups \(\mathcal S_\ell(F,e)\), \(1\le e<t\), have finite
length, and
\begin{equation}\label{eq:host-layer-premise}
 \sum_{e=1}^{t-1}\operatorname{length}_{\mathbb Z_\ell}
                 \mathcal S_\ell(F,e)\le B.
\end{equation}
Here \emph{cyclotomic \(\ell\)-power extension} means
\(F=\mathbb Q(\mu_{\ell m})(\mu_{\ell^b})\) for some \(b\), as
encoded by \texttt{IsLayer}; the statement is invariant under the
number-field presentations used in the source.

Write
\[
 \mathcal S_\ell(E_h,e)^{+_e}
 =\{x\in\mathcal S_\ell(E_h,e):c_{E_h}x=(-1)^e x\},
\]
where the complex-conjugation action includes the Tate coefficients.
Then these eigenspaces have finite length for every odd prime \(\ell\)
and every \(1\le e<t\).  For every finite set \(L\) of odd primes,
\begin{equation}\label{eq:host-point-bound}
 \sum_{\ell\in L}\log\ell\sum_{e=1}^{t-1}
 \operatorname{length}_{\mathbb Z_\ell}
       \mathcal S_\ell(E_h,e)^{+_e}
 \le (Nt)^C+CB.
\end{equation}
There is no Case~I restriction on \(L\) in this assertion.  This is
the precise field-level, finite-parameter statement defined by
\texttt{PointHV}, now with its source inputs supplied.

\paragraph{Closed source inputs and proof.}
Set \(h_N=\texttt{SmallClassNumbers.litN}\), and obtain
\(h_H\) by applying
\path{published_hypotheses_from_patching_alt} to \(h_N\) and the
closed source theorem \texttt{descWeak}.  The following checked
assembly supplies every argument of \texttt{pointHV\_of\_counts}.
The directory prefix in the next list is
\path{ModularityArgument/ThresholdArguments/DenseThresholdCases/}.
\begin{enumerate}
 \item \path{HVFinalU.lean}, \texttt{pt\_s}, gives Poitou--Tate
 duality from \(h_N\); the local compact finiteness and duality
 arguments are the fields \texttt{hN.coh.lfK} and
 \texttt{hN.coh.ltK}.
 \path{KFinal.lean}, \texttt{lemmaK\_litN\_alt}, gives
 \(\exists C_K,\ \texttt{LemmaK}(C_K)\).
 \item \texttt{CountI} is supplied by the closed composition proved
 above.  \texttt{CountII} and \texttt{CountIII} are supplied by
 \path{AdjointLiftingOfCohomology.lean},
 \path{countII_and_countIII_of_printsH_of_sp4In}, using
 \path{Sp4In0.lean}, \texttt{sp4In0\_alt}, at \(h_H\).
 \item \path{DnEndsTop.lean}, \texttt{ends0\_alt}, supplies the
 reduced unitary lifting statement at \(h_H\).  Applying
 \path{LinkedUnlE.lean},
 \path{count_v_of_printed_bound_k_reduced_unitary_lift}, with
 \(h_H\) and the preceding \texttt{LemmaK} witness gives
 \texttt{CountV}.
 \item Outside that directory, \path{LevelCohomology/PropZ.lean}
 proves \texttt{X.propZ} without additional arguments.
 \texttt{countZ X.propZ} gives \texttt{CountZ}.
\end{enumerate}
The construction of \(h_H\) is in
\path{PublishedHypothesesFromPatching.lean}; the exact same closed
\(h_N,h_H\) are used by \texttt{SmallClassNumbers.corGK}.
The theorem \texttt{pointHV\_of\_counts} is in
\path{HeightAndVanishing/JointWinOfCounts.lean}.  Applying it to the
inputs above proves \texttt{PointHV}.  This is a forward proof
composition; no converse implication from the later \texttt{VWin}
theorem is being used.

For clarity, the final arithmetic combination can also be checked
directly.  The source's \texttt{sumOut} bounds the contribution from
odd primes outside \(\mathcal B_\chi\) by \((Nt)^{C_0}\).
For each \(\ell\in\mathcal B_\chi\), its \texttt{sumIn}, using
\eqref{eq:host-layer-premise}, bounds that prime's contribution by
\((Nt)^{C_\ell}+C_\ell B\).  All these bounds include the requisite
finiteness assertions.  Set
\[
 D=\sum_{\ell\in\mathcal B_\chi}C_\ell,
 \qquad b=|\mathcal B_\chi|,
 \qquad C=C_0+D+b+bD.
\]
For \(x=Nt\ge2\), the total is at most
\[
 x^{C_0}+b x^D+bDB
 \le (b+1)x^{C_0+D}+CB
 \le x^{C_0+D+b}+CB
 \le x^C+CB,
\]
because \(b+1\le2^b\le x^b\).  This proves
\eqref{eq:host-point-bound}.  The source proves these steps in
\path{HeightAndVanishing/FinalLe.lean}, \texttt{sumIn};
\path{HeightAndVanishing/SumOut.lean}, \texttt{sumOut}; and
\path{HeightAndVanishing/PointArith.lean}, \texttt{point}.
The layer condition is \texttt{LayerLe} in
\path{HeightAndVanishing/LayerLe.lean};
\texttt{Lchi} in \path{PrimeSetBoundDataOfPrimitiveChar.lean}
is the full cyclotomic field of level \(\ell m\), and
\texttt{IsLayer} is in \path{StmtFields.lean}.

Every assembly declaration just cited, and the complete proofs of
\texttt{sumIn}, \texttt{sumOut}, \texttt{countIV}, and \texttt{point},
were read for this source reconstruction.  The imported arithmetic
and automorphic theorems are used under the same stipulated source
contract as \texttt{CorGK}; this is not a fresh build of the full
Lean project.  The explicit proof composition is recorded separately
as an uncompiled review fragment.

% K17 supporting section, 2026-11-21 UTC.
% Authoring/check model: GPT-6 Astra Pro; effort: not exposed.
% Bibliography entries are supplied separately in k_groups_bibitems.tex.
% Uses the parent's theorem/sourcefact/proof environments and amsmath/amssymb.

\section{The precise consequences for \texorpdfstring{$K$}{K}-groups}
\label{sec:k-groups}

\subsection{The known comparison with Vandiver}

Let \(\ell\) be an odd prime, put
\[
 S_\ell=\operatorname{Spec}\mathbb Z[1/\ell],\qquad
 A_\ell=\operatorname{Cl}(\mathbb Z[\mu_\ell])\otimes\mathbb Z_\ell,
\]
and let \(A_\ell^{[a]}\) denote the eigenspace for the \(a\)-th power
of the Teichm\"uller character, with \(a\) taken modulo \(\ell-1\).
For a finite abelian group \(M\), write \(M\{\ell\}\) for its
\(\ell\)-primary subgroup.

\begin{sourcefact}[Known comparison: BG-KUR-15, BG-WEI-82, BG-WEI-109]
For every \(k\ge1\),
\[
 K_{4k}(\mathbb Z)\{\ell\}
 \simeq H^2(S_\ell,\mathbb Z_\ell(2k+1)),                 \tag{K1}
\]
and
\[
 K_{4k}(\mathbb Z)\{\ell\}=0
 \quad\Longleftrightarrow\quad A_\ell^{[-2k]}=0.         \tag{K2}
\]
The two groups in (K2) are cyclic simultaneously.
Consequently Vandiver's conjecture at \(\ell\) is equivalent to either
of the following conditions:
\[
\begin{split}
 &K_{4k}(\mathbb Z)\{\ell\}=0
       \quad(1\le k\le(\ell-3)/2);\\
 &K_{4k}(\mathbb Z)\{\ell\}=0
       \quad\text{for every }k\ge1.
\end{split}
\tag{K3}
\]
\end{sourcefact}

Here (K1) uses the modern integral comparison edge map, followed by
localization; it is not an unspecified factorial-normalized Chern
character. Exact locators are Weibel \cite[VI.8.2, VI.8.3,
VI.10.9]{Weibel2013} and Kurihara
\cite[Lemma~1.2, Corollary~1.5, Proposition~3.7, Remark~3.11]{Kurihara1992}.
Historically, Kurihara's Proposition~3.7 gave one direction for
\(K\)-groups; his Remark~3.11 made the converse conditional on the
then-conjectural comparison. The modern unconditional equivalence
should therefore not be attributed to an unconditional 1992
comparison theorem.

We record the short character check. Over
\(\mathbb Z[\mu_\ell,1/\ell]\) the omitted prime is principal and the
Brauer group is zero. Kummer theory identifies \(H^2(\mu_\ell)\)
with the class group modulo \(\ell\). Tensoring with
\(\mu_\ell^{\otimes2k}\), and taking invariants under the group of
order \(\ell-1\), selects the exponent \(-2k\). Thus the dimensions
modulo \(\ell\) agree. Finiteness and Nakayama give both the zero
and cyclicity assertions. This argument does not identify the full
finite groups or their orders.

The trivial eigenspace is zero: extension after norm to
\(\mathbb Q\) acts on it by \(\ell-1\), while
\(\operatorname{Cl}(\mathbb Z)=0\). The other even characters are
exactly the exponents \(-2k\) for
\(1\le k\le(\ell-3)/2\). This proves (K3), including the empty range
at \(\ell=3\). For a fixed \(k\), (K2) concerns only one component.
The periodic assertion is about \emph{vanishing}, with period
\((\ell-1)/2\) in \(k\); it is not periodicity of the full finite
groups. In particular \(k=0\), where \(K_0(\mathbb Z)=\mathbb Z\),
is excluded. As a useful coverage consequence, vanishing of
\(K_{4k}\) for \(1\le k\le R\) would imply Vandiver for
every odd \(\ell\le2R+3\) (VK-COVER).

\paragraph{Background cases that do not use K17-S17.}
All regular odd primes have zero primary part in every positive
degree \(4k\). The 2-primary part is also zero in every such degree:
specialize Rognes--Weibel \cite[Theorem~0.7, proved as
Theorem~7.11]{RognesWeibel2000} to \(\mathbb Z[1/2]\), where
\(s=1,t=u=0\), and use localization and finiteness.
Furthermore,
\[
 K_4(\mathbb Z)=0,\qquad K_8(\mathbb Z)=0.              \tag{K4}
\]
The first is Rognes \cite[Theorem~0.1, p.~268]{Rognes2000};
the second is Dutour Sikiri\'c--Elbaz-Vincent--Kupers--Martinet
\cite[Theorem~6.4]{DEKM2026}, published in the January 2026 issue.
An older table listing \(K_8\) as unknown is obsolete.
Hart--Harvey--Ong \cite[\S1, item~(2)]{HartHarveyOng2017}
verify Vandiver for every odd prime \(\ell<2^{31}\).
These facts exclude all such prime divisors of every
\(|K_{4k}(\mathbb Z)|\). Also (K2) and (K4) give vanishing whenever
\(k\equiv0,1,2\pmod{(\ell-1)/2}\); the zero residue uses the
trivial class-group component.

\subsection{The negative Selmer group is exactly the required dual}

Write
\[
 D_\ell(e)=\mathbb Q_\ell/\mathbb Z_\ell(-e),\qquad
 \mathcal S_\ell(e)=
 \ker\!\left(H^1(S_\ell,D_\ell(e))
              \longrightarrow H^1(\mathbb Q_\ell,D_\ell(e))\right).
\]
All cohomology is continuous. The notation agrees with the
rational-field negative Selmer group in K17-S17-Q.

\begin{theorem}[VK-PT and VK-LOCAL]
For odd \(\ell\) and \(k\ge1\), there is a canonical isomorphism
\[
 \bigl(K_{4k}(\mathbb Z)\{\ell\}\bigr)^\vee
       \simeq\mathcal S_\ell(2k).                     \tag{K5}
\]
With the bottom coefficient ring \(O=\mathbb Z_\ell\), the
source's group
\(\operatorname{ShaU}(\chi_\ell^{-2k};2\ell)\)
is this same group.
\end{theorem}

\begin{proof}
Put \(r=2k+1\), \(T=\mathbb Z_\ell(r)\), and
\(D=T^\vee(1)=\mathbb Q_\ell/\mathbb Z_\ell(1-r)\).
At finite level set
\(A_n=T/\ell^nT=\mathbb Z/\ell^n(r)\), with reduction
transition maps, and let
\(D_n=A_n^\vee(1)=\mathbb Z/\ell^n(1-r)\), with the
dual inclusion maps into \(D_{n+1}\).
The finite Poitou--Tate sequence and local duality
\cite[I, Theorem~4.10; II, Theorem~2.9]{MilneADT}
give the segment
\[
\begin{split}
 H^1(\mathbb Q_\ell,A_n)&\longrightarrow H^1(S_\ell,D_n)^\vee
 \longrightarrow H^2(S_\ell,A_n)\\
 &\longrightarrow H^2(\mathbb Q_\ell,A_n)
 \longrightarrow H^0(S_\ell,D_n)^\vee\longrightarrow0.
\end{split}
\tag{K6}
\]
The modified real-place terms vanish since \(\ell\) is odd.
Local duality identifies \(H^2(\mathbb Q_\ell,A_n)\) with
\(H^0(\mathbb Q_\ell,D_n)^\vee\), and the last map is dual to
restriction of invariants. The global and \(\ell\)-local
cyclotomic images are both \(\mathbb Z_\ell^\times\). Their
invariants in \(D_n\) are therefore equal, including when they are
nonzero. Hence the last map in (K6) is an isomorphism.
Exactness and duality therefore give compatible exact sequences
\[
 0\longrightarrow H^2(S_\ell,A_n)^\vee
 \longrightarrow H^1(S_\ell,D_n)
 \longrightarrow H^1(\mathbb Q_\ell,D_n).
\]
This checks the direction of the dual injection before passing
to a limit.

The finite-coefficient arithmetic cohomology groups are finite,
so their inverse systems satisfy the Mittag--Leffler condition.
In particular Jannsen's exact sequence
\cite[(3.1)]{Jannsen1988}
\[
 0\longrightarrow\varprojlim\nolimits^1 H^1(S_\ell,A_n)
 \longrightarrow H^2_{\rm cont}(S_\ell,T)
 \longrightarrow\varprojlim H^2(S_\ell,A_n)
 \longrightarrow0
\]
has zero left term. Pontryagin duality identifies the dual of
this inverse limit of finite groups with
\(\varinjlim H^2(S_\ell,A_n)^\vee\).
Filtered direct limits are exact, and continuous cohomology
with discrete torsion coefficients commutes with this filtered
union: both globally and locally,
\(\varinjlim H^1(D_n)=H^1(D)\).
Taking the direct limit of the finite exact sequences proves
\[
 H^2_{\rm cont}(S_\ell,T)^\vee
 \simeq\ker\!\left(H^1(S_\ell,D)
                   \longrightarrow H^1(\mathbb Q_\ell,D)\right).
\]
Now (K1) gives (K5). Thus neither a local \(H^2\) correction
nor an unexamined inverse-limit term remains.

It remains to compare the explicit source local conditions.
For \(q\ne\ell\), \(D_\ell(e)\) is unramified and
\[
 H^1_{\rm unr}(\mathbb Q_q,D_\ell(e))
 =D_\ell(e)/(q^{-e}-1)D_\ell(e)=0 \qquad(e>0).         \tag{K7}
\]
Indeed \(q^{-e}-1\) is a nonzero \(\ell\)-adic integer and
multiplication by it is onto the divisible group
\(\mathbb Q_\ell/\mathbb Z_\ell\).
The geometric Frobenius convention replaces it by \(q^e-1\)
and gives the same result. Thus local triviality at 2 is
automatic for a class unramified away from \(\ell\).
Local triviality at \(\ell\) implies the inertia condition there.
Conversely, the source's inertia conditions give an
unramified-away-\(\ell\) class, and its decomposition condition
at \(\ell\) makes that class lie in the displayed kernel.
The odd-primary real conditions add nothing.

Finally the recovered definitions give
\(Q_O=O\otimes_{\mathbb Z_\ell}
(\mathbb Q_\ell/\mathbb Z_\ell)\) and
\(e(O)=\operatorname{length}_O(O/\ell O)\).
The bottom ring supplied by \texttt{exists\_coef} is isomorphic
to \(\mathbb Z_\ell\), so \(Q_O=D_\ell(0)\) and \(e(O)=1\).
The scalar twist is exactly \(\chi_\ell^{-e}\).
This proves the asserted normalization and group identification.
\end{proof}

\subsection{Removing every excluded prime}

\begin{theorem}[VK-EXCEPTIONS]
If \(e>0\) is even and \(\ell\) is odd with
\((\ell-1)\mid2e\), then
\[
 K_{2e}(\mathbb Z)\{\ell\}=0.                         \tag{K8}
\]
\end{theorem}

\begin{proof}
The character \(\omega^{-e}\) has order at most two.
The trivial component is already zero. Otherwise
\(-e\equiv(\ell-1)/2\pmod{\ell-1}\).
Evenness of \(e\) forces \(\ell\equiv1\pmod4\).
Put \(L=\mathbb Q(\mu_\ell)\),
\(F=\mathbb Q(\sqrt\ell)\), and \(H=\operatorname{Gal}(L/F)\).
For \(x\in A_\ell^{[-e]}\), the norm and ideal-extension maps
satisfy
\[
 jN(x)=\sum_{h\in H}h(x)=\frac{\ell-1}{2}x.           \tag{K9}
\]

We need only the elementary estimate \(h_F<\ell\).
Put \(M=\lfloor\sqrt\ell/2\rfloor\).
Minkowski's bound \cite[Theorem~4.3, p.~70]{MilneANT}
provides an integral ideal of norm at most \(M\) in each class.
If \(a_F(n)\) is the number of ideals of norm \(n\), then
\(a_F(n)\le d(n)\): at a split prime the coefficient of
the \(a\)-th power is \(a+1\), at an inert prime it is zero or
one, and at a ramified prime it is one. Multiplicativity gives
the claim. Hence
\[
 h_F\le\sum_{n=1}^M a_F(n)
 \le\sum_{n=1}^M d(n)
 =\#\{(a,b)\in\mathbb Z_{>0}^2:ab\le M\}
 \le M^2\le\frac{\ell}{4}<\ell.                      \tag{K10}
\]
Here \(\ell\ge5\) in the nontrivial quadratic case.
Thus \(\operatorname{Cl}(F)\{\ell\}=0\), so \(N(x)=0\).
The scalar \((\ell-1)/2\) is a unit on an \(\ell\)-primary
group; (K9) gives \(x=0\). Apply (K2) with \(k=e/2\).
No integral injectivity of an ideal-extension map is assumed.
\end{proof}

\subsection{The global order and cyclicity estimates}

\begin{theorem}[VK-POLY and VK-CYCLIC-RANGE]
The recovered theorem K17-S17-Q implies absolute
\(C,A\in\mathbb N\) such that, for every \(k\ge1\),
\[
 \log|K_{4k}(\mathbb Z)|\le C(2k+1)^A.               \tag{K11}
\]
Put \(B_k=C(2k+1)^A\). For every prime \(\ell\),
\[
 v_\ell\bigl(|K_{4k}(\mathbb Z)|\bigr)
       \le\left\lfloor\frac{B_k}{\log\ell}\right\rfloor. \tag{K12}
\]
In particular \(\log\ell>B_k\) implies vanishing of the primary
part, while \(2\log\ell>B_k\) makes it either zero or
\(\mathbb Z/\ell\).
\end{theorem}

\begin{proof}
The positive even \(K\)-group is finite. Its 2-primary part is
zero by the cited theorem; the 3-primary part is zero by the
trivial-character case of (K2). All primes excluded by
\((\ell-1)\mid4k\) vanish by (K8).
Apply K17-S17-Q to the finite set of remaining prime divisors
of its order. By (K5), the normalized source lengths are
exactly
\(v_\ell(|K_{4k}(\mathbb Z)|)\). The finite weighted sum is
therefore the logarithm of the order, proving (K11).
Singleton estimates give (K12). A finite
\(\mathbb Z_\ell\)-module of length at most one is zero or
\(\mathbb Z/\ell\), proving the last assertion.
\end{proof}

For odd \(\ell\), Kurihara's cyclicity comparison also makes
\(A_\ell^{[-2k]}\) cyclic in the last range. It does not
bound that class-group component's order by \(\ell\);
the comparison controls its number of generators, not its
whole order.

\subsection{Extracting exponent seven from the source proof}

\begin{theorem}[VK-EXP7]
Keeping the source's fixed levels and using the exact
highest-weight dimensions, the bound can be sharpened to
\[
 \exists C_7\in\mathbb N\quad
 \log|K_{4k}(\mathbb Z)|\le C_7(2k)^7
 \quad\text{for every }k\ge1.                       \tag{K13}
\]
The constant has not been numerically evaluated.
\end{theorem}

\begin{proof}
We identify the representation-theoretic input and then give
the full count-to-norm argument. All automorphy, cohomology,
and representation-classification inputs remain the cited
source theorems used in K17-S17.

In \path{WeylC.lean}, indices start at zero and
\[
 \nu_j=j+1+H-m_{n-1-j},\qquad N_{\rm det}=n+H.
\]
The helper \path{ExistsScalar.lean:existsScalar} calls
\path{Scalar.lean:scalar}, which constructs the polynomial
module generated by a nonzero highest-weight vector of
weight \(\nu\). This module is precisely one copy of
the irreducible \(L_\nu\). To check this assertion, decompose
the finite polynomial representation into irreducibles.
Every nonzero component of the highest-weight vector has
highest weight \(\nu\); in the isotypic component
\(L_\nu\otimes W\), that vector is a highest-weight vector
in \(L_\nu\) tensored with one vector of \(W\). Its orbit
span is one copy of \(L_\nu\).
The implementation takes the span over all matrices.
It equals the invertible-matrix orbit span: any linear
functional vanishing on the latter gives a polynomial that
vanishes on the Zariski-dense invertible locus, hence on all
matrices. Thus this exact dimension applies to the actual
constructed module, not just an alternative informal choice.

The helper \path{ExistsTens.lean:exists_tens} gives the
required irreducible coefficient representation \(\Xi\)
after determinant twist \(-N_{\rm det}\), with
\[
 \dim\Xi\le(\dim L_\nu)^{[F:\mathbb Q]}
\]
for the totally complex fields used here. The infinitesimal
character is the same one used by the source's cohomology
theorem. The standard highest-weight classification and Weyl
dimension formula are cited background
\cite[\S\S27.2--27.3 and \S28.1]{EtingofGL2024}:
\[
 \dim L_\nu=\prod_{0\le i<j<n}
       \frac{\nu_i-\nu_j+j-i}{j-i}.                  \tag{K14}
\]

The rank-three trace-count source uses
\(m=(e-1,-1,-e-1)\), \(H=e+1\).
It gives
\[
 \nu=(2e+3,e+4,5),\qquad
 \lambda=\nu-(e+4)(1,1,1)=(e-1,0,1-e).
\]
Formula (K14) is \(e^3\), and an imaginary quadratic
field has two complex embeddings, so \(\dim\Xi_3\le e^6\).
The rank-two branch uses \(m=(e,-e)\), \(H=e\), giving
\(\nu=(2e+1,2)\) and \(\lambda=(e-1,-e)\).
Its one-embedding dimension is \(2e\); hence
\(\dim\Xi_2\le4e^2\).

Fix \(K=\mathbb Q(\sqrt{-3})\), \(N=2\).
The proof of \texttt{cardK} counts pairwise nonisomorphic
cuspidal representations by total cohomology at
\(\operatorname{KLev}(K,3,6)\).
The theorem \path{CohBoundFree.lean:cohBoundFree} therefore
gives a fixed positive integer \(B_0\) such that this total
cohomology is at most \(B_0\dim\Xi_3\).
It follows, by the unchanged
\path{CountSys.lean:exists_eigData} construction, that
the rank-three trace data have at most \(B_0e^6\) elements.

For the rational rank-two branch the level is
\(N|\operatorname{disc}K|=6\).
The source's quadratic-base-change argument uses six fixed
imaginary quadratic fields with pairwise coprime discriminants.
For concreteness one may take discriminants
\(-3,-7,-11,-19,-23,-31\).
At any prime \(p\), five of them avoid \(p\).
For each field choose a fixed positive integer \(B_i\)
bounding total cohomology at
\(\operatorname{KLev}(K_i,2,18)\) by \(B_i\dim\Xi_2\).
These seven cohomology constants are independent of
\(e,p,\ell\).

The proof of
\path{CountQuadraticOfHyp.lean:count_quadratic_of_hypotheses}
then applies unchanged. Its helper
\path{RootSum.lean:rootSum} introduces at most \(2^n\)
choices of roots of Frobenius eigenvalues and a Boolean
factor for two possible lifts. For \(n=2\) the multiplier
is \(8\). Thus, if \(I_p\) is the five-element subset
avoiding \(p\), the rational trace-data cardinality is at most
\[
 8\sum_{i\in I_p}4B_i e^2
 \le32\left(\sum_{i=1}^6B_i\right)e^2.
\]
This indexes whole eigensystems, retaining the multiplicities
needed by the norm product. Replacing them by distinct
single-prime trace values would not justify this argument.
Both trace-conversion helpers preserve indexing cardinality.
With
\[
 D=B_0+32\sum_{i=1}^6B_i,
\]
their disjoint union therefore has at most \(De^6\) elements.

Let \(p\ge7\), \(p\nmid6\), choose \(w\mid p\) in \(K\), and put
\(q=Nw=p^f\), \(1\le f\le2\). For the trivial character
orbit, the source constructs the integer
\[
 n_p(e)=\prod_{s}
       q^e\bigl(t_s-(q^{-e}+1+q^e)\bigr)
       \in\mathbb Z\setminus\{0\}.                  \tag{K15}
\]
The data are Galois stable, each scaled factor is integral,
and the purity bound \(|t_s|\le3\) makes every factor nonzero.
The exact source theorem is
\path{FrobeniusTraceNormNonzeroLogBound.lean:}
\texttt{frobenius\_trace\_norm\_nonzero\_log\_bound}.
It gives
\[
\begin{split}
 \log|n_p(e)|
 &\le De^6(2e\log q+1)\\
 &\le5De^7\log p,
\end{split}
\tag{K16}
\]
since \(q\le p^2\), \(e\ge1\), and \(\log p\ge1\).
The source's unchanged \texttt{OneChar} and orbit argument
then bounds the normalized negative-Selmer length by
\(2v_\ell(n_p(e))\), for \(\ell\ne p\) in the source's
admissible range. Its other arithmetic hypotheses are
exactly those checked by \path{ClassNumberThree.lean:gd}.

Take \(e=2k\) and \(p=7\). All excluded primary parts are
zero by (K8) and the known 2-primary theorem.
The one additional primary part \(\ell=7\) is zero because
7 is regular. Thus every prime satisfies
\[
 v_\ell(|K_{4k}(\mathbb Z)|)\le2v_\ell(n_7(2k)).
\]
Consequently
\[
 |K_{4k}(\mathbb Z)|\mid|n_7(2k)|^2,\qquad
 \log|K_{4k}(\mathbb Z)|
 \le10D\log7\,(2k)^7<20D(2k)^7.
\]
Taking \(C_7=20D\) proves (K13).
\end{proof}

This is a mathematical sharpening of the recovered proof,
not a declaration that the separate proposition
\path{HeightAndVanishing/Counts.lean:CorG} has been proved.
The exact source proof witnesses, without these improvements,
give exponent \(133\): the rank-three exponent is \(36\),
the six rank-two exponents are \(16\) each, and the norm
estimate adds one. Replacing sums of exponents by maxima
already gives \(37\). The exponent \(7\) additionally uses
the fixed levels and the exact coefficient dimensions.
No claim of a new compiled Lean theorem is made.

\paragraph{An exact small-degree computation target: VK-TRACE-GCD.}
Use (K15) also at \(p=11\). Since the 11-primary part
already vanishes,
\[
 |K_{4k}(\mathbb Z)|
 \mid \gcd\bigl(|n_7(2k)|,|n_{11}(2k)|\bigr)^2.       \tag{K17}
\]
Indeed each valuation is bounded by twice each of the two
integer valuations. Here \(q_7=7\), whereas \(q_{11}=121\).
A complete computation of these trace data and their gcd
could prove \(K_{12}=0\) at \(e=6\), or \(K_{16}=0\)
at \(e=8\), if the gcd has only already-excluded prime factors.
Neither trace computation has been completed in this run.
Formula (K17) identifies the finite calculation that would
be needed, including its normalization and multiplicities.

\subsection{An exact polynomial certificate for a new vanishing degree}

The remaining finite computation can be expressed without choosing
algebraic eigenvalues individually. Retain the entire indexed trace
multiset from the proof of VK-EXP7, including repeated values.

\begin{proposition}[VK-POLY-CERT]
Let $e=2k>0$ and let $p\ge7$ be a regular auxiliary prime. Put
$q=\operatorname{Nm}(w)$ for the chosen place of
$\mathbb Q(\sqrt{-3})$ above $p$, and define
\[
 F_{p,e}(X)=\prod_{s\in S_{p,e}}(X-q^e t_s),\qquad
 E_{p,e}=q^{2e}+q^e+1.
\]
Then $F_{p,e}$ is monic in $\mathbb Z[X]$, its value at $E_{p,e}$
is nonzero, and
\[
 |K_{2e}(\mathbb Z)|\mid |F_{p,e}(E_{p,e})|^2.
 \tag{K18a}
\]
More generally, let $P_{p,e}\in\mathbb Z[X]$ be monic and nonzero,
with $F_{p,e}\mid P_{p,e}$ in $\mathbb Q[X]$. Remove the maximal
power of $X-E_{p,e}$ from $P_{p,e}$, obtaining $P^*_{p,e}$.
Then
\[
 P^*_{p,e}(E_{p,e})\ne0,\qquad
 |K_{2e}(\mathbb Z)|\mid |P^*_{p,e}(E_{p,e})|^2.
 \tag{K18b}
\]
For two such auxiliary primes the order divides the square of the
gcd of these two nonzero evaluations. If every prime factor of this
gcd is below $2^{31}$, the already verified Vandiver range implies
$K_{2e}(\mathbb Z)=0$.
\end{proposition}

\begin{proof}
The source properties \texttt{TrData.int} and
\texttt{TrData.conj} in \path{AutomorphicVarieties/AutomorphicForms/Classifies.lean} say that every
$q^e t_s$ is an algebraic integer and that every Galois automorphism
permutes the indexed multiset. The elementary symmetric
coefficients are therefore rational algebraic integers. This gives
the asserted integral monic polynomial. Its evaluation satisfies
\[
 n_p(e)=(-1)^{|S_{p,e}|}F_{p,e}(E_{p,e}).
\]
The bound $|t_s|\le3$ and $q\ge7$ imply
$E_{p,e}>3q^e$, so every evaluated factor is nonzero.
The all-prime divisibility in VK-EXP7, using regularity to remove
the auxiliary prime itself, gives (K18a).

Monic polynomial division gives $P_{p,e}=F_{p,e}Q$ with
$Q\in\mathbb Z[X]$. Each exact division by the monic linear
polynomial $X-E_{p,e}$ also stays in $\mathbb Z[X]$.
Since $F_{p,e}(E_{p,e})\ne0$, $F_{p,e}$ has no such linear
factor. Unique factorization in $\mathbb Q[X]$ shows that it still
divides $P^*_{p,e}$; monic division makes the quotient integral.
This proves (K18b). Taking valuations at each prime proves the
gcd assertion. The final assertion uses the already cited vanishing
of every primary part at primes below $2^{31}$.
\end{proof}

Thus a proved, multiplicity-preserving inclusion of the source
trace system in a larger integral Hecke model is sufficient.
Removing the exact scalar root above can avoid a separate
Eisenstein projection. Taking the squarefree part is not permitted:
distinct representations may have the same trace at one prime,
and their multiplicities occur in the source integer.

Both $7$ and $13$ split in $\mathbb Q(\sqrt{-3})$ and are regular.
The following exact evaluation points were checked by integer
arithmetic; no corresponding Hecke polynomials were computed.
\[
\begin{array}{c|r|r|r}
K\text{-group}&e&p&E_{p,e}\\\hline
K_{12}&6&7&13\,841\,404\,851\\
K_{12}&6&13&23\,298\,089\,949\,291\\
K_{16}&8&7&33\,232\,936\,334\,403\\
K_{16}&8&13&665\,416\,609\,998\,910\,563
\end{array}
\]
The required highest weights are $(5,0,-5)$ and $(7,0,-7)$.
The open input is a complete integral Hecke model at the source's
specified levels, together with the coverage and normalization
proof. A dimension count or a spectrum computed from a
conjectural trace formula does not supply that input. The
literature and available-software checks, and this precise finite
certificate reduction, are recorded in
\path{evidence/work/k_numeric_followup.md}.

\subsection{The full-domain Soul\'e index}

\begin{theorem}[VK-SOULE and VK-SOULE-NORM]
Let \(\ell\) be odd and \(m\ge3\) odd, and put
\[
 N_\ell=\operatorname{Gal}
   (\overline{\mathbb Q}/\mathbb Q(\mu_{\ell^\infty})).
\]
For the standard Soul\'e character \(\kappa_m\) on this
full cyclotomic kernel,
\[
 [\,\mathbb Z_\ell(m):\kappa_m(N_\ell)\,]
   =|H^2(S_\ell,\mathbb Z_\ell(m))|
   =|K_{2m-2}(\mathbb Z)\{\ell\}|.                   \tag{K18}
\]
\end{theorem}

\begin{proof}
Fix compatible roots \(\zeta_n\) of order \(\ell^n\).
The standard character is given by Kummer action on roots of
\[
 E_{m,n}=
 \prod_{a\in(\mathbb Z/\ell^n)^\times}
       (1-\zeta_n^a)^{a^{m-1}}.                    \tag{K19}
\]
Kurihara \cite[p.~233, (4)]{Kurihara1992} defines the
cyclotomic class \(c(1)\in H^1(S_\ell,\mathbb Z_\ell(m))\)
at finite level by corestriction of
\((1-\zeta_n)\otimes\zeta_n^{\otimes(m-1)}\).
Restriction--corestriction to \(N_\ell\) gives the sum
of its conjugates. Factoring out
\(\zeta_n^{\otimes(m-1)}\) produces precisely (K19).
Thus \(\operatorname{res}_{N_\ell}c(1)=\kappa_m\);
there is no factorial in this normalization.

The group \(H^1(S_\ell,\mathbb Z_\ell(m))\) is free of
rank one and \(c(1)\ne0\), as recalled on the same page.
Let \(\alpha\) be a primitive generator. Its reduction in
\(H^1(S_\ell,\mathbb F_\ell(m))\) is nonzero by the
coefficient exact sequence. Write
\(\Gamma=\mathbb Z_\ell^\times=\Delta\times\Gamma_1\),
where \(|\Delta|=\ell-1\) and \(\Gamma_1=1+\ell\mathbb Z_\ell\).
Since \(m\) is odd, its residual character on \(\Delta\)
is nontrivial. Inflation--restriction gives
\[
 H^1(\Gamma,\mathbb F_\ell(m))
 =H^1(\Gamma_1,\mathbb F_\ell(m))^\Delta=0.           \tag{K20}
\]
Here \(\Gamma\) is abelian, so \(\Delta\) acts trivially
on \(\Gamma_1\) and by the nontrivial character on the
coefficient module.

If \(\alpha|_{N_\ell}\) had image contained in
\(\ell\mathbb Z_\ell(m)\), its reduction would restrict
to zero. Inflation--restriction for the cyclotomic quotient
would put that nonzero reduction in the zero group (K20),
a contradiction. Apply this first in the maximal extension
unramified outside \(\ell\); the full group \(N_\ell\)
surjects onto its cyclotomic kernel. The image of
\(\alpha|_{N_\ell}\) is a compact, hence closed subgroup
of \(\mathbb Z_\ell(m)\). Since it contains a unit,
it is the whole module.
Writing \(c(1)=a\alpha\), \(a\ne0\), now gives
\[
 [\mathbb Z_\ell(m):\kappa_m(N_\ell)]
 =[H^1(S_\ell,\mathbb Z_\ell(m)):\mathbb Z_\ell c(1)].
\]
Kurihara's Proposition~5.1 gives an inequality; his
Remark~5.2 explicitly states equality, citing Bloch--Kato,
\S6, (6.8)--(6.9). We use that equality as a known
theorem. It identifies the last index with \(|H^2|\).
The modern comparison (K1) completes (K18).
\end{proof}

This reconciles the index formula with Ishii's precise
full-domain statement \cite[Theorem~1.1, p.~2]{IshiiKummer2024}:
every odd-degree character is nonzero; the character is
surjective when \(m\equiv1\pmod{\ell-1}\); and all odd-degree
characters are surjective exactly when Vandiver holds at
\(\ell\). For a fixed degree the condition is the one
component in (K2).

Changing the chosen radical does not change the character
on \(N_\ell\), which fixes the roots of unity.
Replacing \(1-\zeta_n^a\) by \(\zeta_n^a-1\) changes a
Kummer unit by a sign, whose odd-primary Kummer class is
zero. Changing the compatible root generator by
\(b\in\mathbb Z_\ell^\times\) changes the scalar coordinate
by the unit \(b^{-m}\), so the index is unchanged.
Half-products or real norms may introduce powers of 2,
also units here. Ichimura--Kaneko's universal-series
coordinate \cite[\S1.1, pp.~313--314]{IchimuraKaneko1989}
contains \((1-\ell^{m-1})^{-1}\kappa_m/m!\):
the first factor is a unit, whereas \(m!\) need not be.
It must not be inserted into the definition (K19).
The original Ichimura--Sakaguchi paper was bibliographically
identified but not directly read in this run; no original
theorem number from it is asserted.

The proof uses oddness of \(\ell\), especially (K20), and
does not extend the index formula to 2. Nor does it
identify the image of \(\kappa_m\) on a deeper Galois
filtration subgroup with its image on all of \(N_\ell\).
Consequently it proves no integral free-generation theorem.

\subsection{A reproducible certificate at 37}

The criterion of Buhler--Harvey
\cite[\S4, author preprint pp.~4--5]{BuhlerHarvey2011}
says to test, for every irregular index \(a\) at \(\ell\),
\[
 V_{\ell,a}=
 \prod_{c=1}^{(\ell-1)/2}
   (z^c-z^{-c})^{c^{\ell-1-a}}\pmod q,\qquad
 V_{\ell,a}^{(q-1)/\ell}\not\equiv1\pmod q,           \tag{K21}
\]
where \(q\equiv1\pmod\ell\) is prime and \(z\) has order
\(\ell\) modulo \(q\).
For \(\ell=37\), the residues of
\(B_2,B_4,\ldots,B_{34}\) modulo 37 are
\[
31,16,15,16,4,17,32,22,31,15,15,17,12,29,2,0,2.
\]
Thus the only irregular index is \(a=32\).
Taking \(q=149\), \(z=16\), the exact computation gives
\[
 16^{37}\equiv1,\quad16\not\equiv1,\quad
 V_{37,32}=48,\quad48^4=142\not\equiv1\pmod{149}.
\]
Trial division certifies the primes and hence the root's
exact order. This proves Vandiver at 37 by the cited
criterion, and consequently
\[
 K_{4k}(\mathbb Z)\{37\}=0\quad(k\ge1).
\]
The accompanying exact-arithmetic script compares two
Bernoulli algorithms and two modular-product algorithms.
The code and the application of the criterion were read
in full after running. This is a reproducible numerical
certificate with a cited arithmetic criterion, not a new
formal proof certificate (VK-37).

\subsection{The remaining numerical question}

The known primary-part exclusions give the sufficient gate
\[
 |K_{4k}(\mathbb Z)|<2^{31}
 \quad\Longrightarrow\quad K_{4k}(\mathbb Z)=0.       \tag{K22}
\]
A nontrivial finite group has a prime divisor of its order,
and all smaller primes have been excluded. Thus either a
small enough explicit bound or the exact trace gcd in
(K17) could close a fixed-degree case. Merely knowing that
an absolute constant in (K11) or (K13) exists does not.
Indeed the source's particular coarse choice
\(C=36c_5\,2^A\) satisfies \(C\ge36>\log(2^{31})\)
whenever \(c_5\ge1\); that bound cannot meet the gate even
with its smallest positive integer constant.

For a precise logical illustration, choose a prime
\(q>2^{31}\), put \(h=(q-1)/2\), and define abstract groups
\[
 M_k=
 \begin{cases}
 \mathbb Z/q,& k\equiv3\pmod h,\\
 0,&\text{otherwise}.
 \end{cases}
\]
They have no 2-primary or below-\(2^{31}\) primary parts,
\(M_1=M_2=0\), the prescribed fixed-prime vanishing
periodicity, and
\(\log|M_k|\le\lceil\log q\rceil(2k+1)^7\), but
\(M_3\ne0\). These are only a model of the specifically
listed size and numerical patterns. They are not actual
\(K\)-groups or cyclotomic class data and give no
counterexample to Vandiver (VK-NUM-GATE).

The current primary-source review found no proof of
\(K_{12}(\mathbb Z)=0\) or \(K_{16}(\mathbb Z)=0\).
This reports the scope of the search, not an exhaustive
negative theorem about the literature.
Kupers \cite[Remark~2.6, p.~4]{Kupers2019} gives an
integral Church--Farb--Putman route to \(K_{12}\).
The subsequent result of Br\"uck--Miller--Patzt--Sroka--Wilson
\cite[Theorem~B]{BrueckEtAl2024} proves rational
codimension-two vanishing, and does not state the integral
and higher-codimension input required by that route.
No new specified degree vanishing is claimed from the
existential constants extracted here.

\section{Tripod kernels and integral Soul\'e generators}
\label{sec:tripod-integral}

\subsection{Correct field formulation and the odd regular theorem}

For a prime $p$, write
\[
 X=\mathbb P^1_{\mathbb Q}\setminus\{0,1,\infty\},\qquad
 \rho_p:G_{\mathbb Q}\longrightarrow
 \operatorname{Out}\bigl(\pi_1^{\mathrm{et}}(X_{\overline{\mathbb Q}})^{(p)}\bigr),
\]
\[
 F=\mathbb Q(\mu_p),\quad K=\mathbb Q(\mu_{p^\infty}),\quad
 \Omega_p^*=\overline{\mathbb Q}^{\ker\rho_p},
\]
and let $\Omega_p$ be the maximal pro-$p$ extension of $F$ unramified outside $p$. Thus the desired equality is
\[
 \Omega_p^*=\Omega_p,
 \quad\text{equivalently}\quad
 \ker\rho_p=\operatorname{Gal}(\overline{\mathbb Q}/\Omega_p).
\]
The kernel is a subgroup of $G_{\mathbb Q}$; it is not itself the maximal ramification Galois group. At $p=2$ the ramification convention at infinity is specified in Theorem~\ref{thm:pro2-kernel}.

Put $G_p=\operatorname{Gal}(\Omega_p^*/K)$, with the Ihara weight filtration $F^mG_p$, and
\[
 \mathfrak g_p=\bigoplus_{m\geq1}F^mG_p/F^{m+1}G_p.
\]
For odd $m\geq3$, choose $\sigma_m\in F^mG_p$ whose Soul\'e value generates the actual image ideal $\kappa_m(F^mG_p)$. Let
\[
 \psi:\operatorname{Lie}_{\mathbb Z_p}(s_3,s_5,\ldots)
       \longrightarrow\mathfrak g_p,
 \qquad s_m\longmapsto\sigma_m,
\]
and let $\Psi$ be the corresponding homomorphism from the free pro-$p$ group on these generators.

\textbf{Sharifi's exact theorem.} For an \emph{odd regular} prime $p$, $\Psi$ is onto. If rational Deligne--Ihara holds, then $\psi$ and $\Psi$ are isomorphisms and $\Omega_p^*=\Omega_p$ \cite[Theorem~1.1, author's preprint p.~2]{Sharifi2002}. Brown supplies the rational conjecture \cite{Brown2012,Brown2019}. Therefore both the kernel-field equality and integral odd free generation hold for odd regular primes. This resolves the correctly formulated known row LM6~A.44 without step~17. Ishii's recent primary account restates exactly this scope \cite[Theorem~1.1 and Remark~1.1.1, p.~2]{IshiiKernel2026}.

\subsection{A separate deduction for the prime \texorpdfstring{$2$}{2}}
\label{sec:pro2-proof}

The following result is a deduction developed and checked in this run. We did not locate it as a stated theorem in the literature searched, and we make no novelty claim. Its proof uses established arithmetic and rational Deligne--Ihara results together with the complete argument below. It does not use the proposed step--17 estimate. The initiating researcher, two other agents, and the parent researcher each read the whole argument; the check ledger distinguishes those reviews from the cited literature inputs and from formal verification.

\begin{theorem}[K17-P2-MOMENT]\label{thm:pro2-kernel}
Let $X=\mathbb P^1_{\mathbb Q}\setminus\{0,1,\infty\}$ and $\Pi=\pi_1^{\mathrm{et}}(X_{\overline{\mathbb Q}})^{(2)}$. Let $\Omega_2$ be the maximal pro-$2$ extension of $\mathbb Q$ unramified at every finite prime other than $2$, with complexification at the real place allowed. If $\Omega_2^*$ is the fixed field of the kernel of the outer representation
\[
 \rho_2:G_{\mathbb Q}\longrightarrow\operatorname{Out}(\Pi),
\]
then $\Omega_2^*=\Omega_2$.
\end{theorem}

The convention at infinity is part of the statement. The theorem concerns faithfulness of a pro-$2$ action; it does not assert that its integral weight Lie algebra is freely generated by odd Soul\'e elements.

\paragraph{Factorization through the maximal ramification group.}
The image of $\rho_2$ is pro-$2$. Here is a direct justification that does not invoke an odd-prime formulation. The geometric group $\Pi$ is free pro-$2$ of rank two, and its Frattini quotient is $\mathbb F_2(1)^2$. The Galois action on this quotient is trivial. For any finite $2$-group $Q$, the kernel of
\[
 \operatorname{Aut}(Q)\longrightarrow
 \operatorname{Aut}(Q/\Phi(Q))
\]
is a $2$-group. To see this directly, fix a basis of $Q/\Phi(Q)$ of size $d$. Its lifts form a set of cardinality $|\Phi(Q)|^d$, and every such tuple generates $Q$ by the defining property of the Frattini subgroup. The displayed kernel acts freely on this set, since an automorphism fixing a generating tuple is the identity. Its order therefore divides a power of $2$. Use the characteristic finite quotients obtained from the descending Frattini series of $\Pi$. They are cofinal among its finite quotients, and an automorphism trivial on $\Pi/\Phi(\Pi)$ induces an automorphism trivial on the Frattini quotient of each of them. Consequently the corresponding kernel in $\operatorname{Aut}(\Pi)$ is pro-$2$. Inner automorphisms lie in that kernel, so the kernel in $\operatorname{Out}(\Pi)$ is also pro-$2$.

At every odd prime, the marked curve $(\mathbb P^1,\{0,1,\infty\})$ has good reduction with disjoint marked sections. Prime-to-residue-characteristic specialization therefore makes inertia act trivially on its outer pro-$2$ fundamental group. Equivalently, one may use the all-primes factorization recalled in Hain--Matsumoto, \S8.4, p.~27 \cite{HainMatsumoto2001}. Combining unramifiedness outside $2$ with the pro-$2$ image gives an induced continuous representation of
\[
 P=\operatorname{Gal}(\Omega_2/\mathbb Q).
\]
It suffices to prove that representation faithful.

\paragraph{The arithmetic source.}
Gras--Jaulent's Theorem/Definition~0(iv) and Theorem~2 give
\[
 P\simeq C_2*\mathbb Z_2
\]
for the unrestricted-at-infinity group \cite[pp.~1--3]{GrasJaulent2008}. The specialization of their $2$-rationality hypothesis to $\mathbb Q$ is elementary. A rational number whose valuations at all odd primes are even is, modulo rational squares, $\epsilon2^e$ with $\epsilon\in\{1,-1\}$ and $e\in\{0,1\}$. Being a square in $\mathbb Q_2$ forces $e=0$ by valuation and $\epsilon=1$ because $-1$ is not a $2$-adic square. The global and local $2$-primary root-of-unity groups both have $2$-rank one. Thus their criterion applies.

Kronecker--Weber identifies the full abelianization with the cyclotomic character
\[
 \chi:P\longrightarrow\mathbb Z_2^\times.
\]
Write $P=\langle c\rangle*\langle\gamma\rangle$, where $c^2=1$ and $\langle\gamma\rangle\simeq\mathbb Z_2$. Then $\chi(c)=-1$. Replacing $\gamma$ by $c\gamma$ if necessary is an automorphism of the free product and ensures that
\[
 u=\chi(\gamma)\in1+4\mathbb Z_2.
\]
Since $\chi$ is the abelianization, $u$ is a topological generator of $1+4\mathbb Z_2$.

The subgroup $P_i=\ker(\chi\bmod4)$, corresponding to $\mathbb Q(i)$, is free pro-$2$ on $\gamma$ and $c\gamma c$. One can verify this without an infinite-rank subgroup theorem: let the involution exchange the two generators $a,b$ of a free pro-$2$ group. Its semidirect product with $C_2$ has the universal property of $C_2*\mathbb Z_2$, with $a=\gamma$ and $b=c\gamma c$. Its index-two kernel is the asserted free group. Set
\[
 h=(c\gamma c)\gamma^{-1},\qquad H=\ker\chi.
\]
Then $(\gamma,h)$ is another free pro-$2$ basis of $P_i$, and $H$ is the kernel of the map $P_i\to\mathbb Z_2$ sending $\gamma$ to $1$ and $h$ to $0$. For $x\in\mathbb Z_2$, write
\[
 h_x=\gamma^xh\gamma^{-x}.
\]

\paragraph{Finite quotients and measure-valued Magnus coefficients.}
Let $N_n$ be the closed normal subgroup of $P_i$ generated by $\gamma^{2^n}$. Then
\[
 P_i/N_n\simeq C_{2^n}*\mathbb Z_2.
\]
The kernel $H_n$ of its map to $C_{2^n}$ is free pro-$2$ on the $2^n$ elements $h_j$, $j\bmod2^n$. Indeed, the free pro-$2$ group on those symbols, with the cyclic permutation action, has semidirect product with $C_{2^n}$ possessing precisely the universal property of $C_{2^n}*\mathbb Z_2$.

The subgroups $N_n$ decrease to $1$. In any finite continuous $2$-group quotient of $P_i$, the image of $\gamma$ has finite $2$-power order, so $N_n$ lies in its kernel for all sufficiently large $n$. Compactness also makes the map to the inverse limit surjective: compatible quotient cosets form a nested family of nonempty closed subsets of $P_i$. Hence
\[
 P_i\simeq\varprojlim_nP_i/N_n,
 \qquad
 H\simeq\varprojlim_nH_n.                                  \tag{P2.1}
\]
The transition maps fold the finite bases by reduction of $j$ modulo $2^n$. The elements $h_x$ topologically generate $H$. To check this last point explicitly, the kernels of $H\to H_n$ decrease to $1$, and compactness implies that every open normal subgroup contains one of them. Every finite quotient of $H$ therefore factors through an $H_n$, whose indicated generators are images of the $h_x$.

The finite-rank integral Magnus maps are injective. A primary locator is \cite[\S5.1, p.~21, (5.1)]{MinacPasiniQuadrelliTan2020}, whose mod-$2$ statement implies the integral injection. There is also a short direct verification. The augmentation ideal of $\mathbb F_2[Q]$ for a finite $2$-group $Q$ is nilpotent: choose a central element $z$ of order two, use $(z-1)^2=0$, and induct on $Q/\langle z\rangle$. Thus any map from a finite-rank free pro-$2$ group to $Q$ is recovered by substituting $X_j\mapsto\overline g_j-1$ into the mod-$2$ Magnus series. An element with trivial expansion maps to $1$ in every finite quotient and is therefore $1$.

The degree-$r$ coefficients of the compatible integral expansions in (P2.1) belong to
\[
 M_r=\varprojlim_n\mathbb Z_2[(\mathbb Z/2^n)^r]
     =\mathbb Z_2[[\mathbb Z_2^r]],
\]
the bounded $\mathbb Z_2$-valued measures on $\mathbb Z_2^r$. With multiplication given by external product and concatenation, they define an injective homomorphism
\[
 \theta:H\longrightarrow1+\prod_{r\geq1}M_r,
 \qquad \theta(h_x)=1+\delta_x.                             \tag{P2.2}
\]
Injectivity follows from the finite-rank maps and (P2.1). Consequently every nontrivial element has a first finite nonzero Magnus degree. Completion in the $(2,I)$ topology creates no element with all homogeneous coefficients zero: the completed coefficient algebra is the displayed product.

\paragraph{The actual rational image.}
Let $B$ be the proalgebraic closure over $\mathbb Q_2$ of the rational outer image induced by the tripod representation, and let $U$ be its prounipotent radical. Hain--Matsumoto, \S8.3 and Proposition~8.3, Theorem~8.4, pp.~25--26, gives an extension
\[
 1\longrightarrow U\longrightarrow B\longrightarrow\mathbb G_m\longrightarrow1
\]
whose reductive quotient is induced by $\chi$, and a continuous map
\[
 \rho_U:H\longrightarrow U(\mathbb Q_2)                    \tag{P2.3}
\]
with Zariski-dense image \cite{HainMatsumoto2001}. The rational action is induced functorially from the pro-$2$ action; a trivial pro-$2$ outer automorphism gives a trivial rational outer automorphism. Their Appendix~A.6, A.9 and A.10 records the relevant comparison. We have used the actual arithmetic image, not an unrelated motivic group.

Brown's rational Deligne--Ihara theorem and Hain--Matsumoto's graded comparison give, with positive arithmetic weights,
\[
 \operatorname{gr}\operatorname{Lie}U
 \simeq\operatorname{Lie}_{\mathbb Q_2}(s_3,s_5,s_7,\ldots),
 \qquad\operatorname{wt}(s_m)=m.                           \tag{P2.4}
\]
The exact outer-action statement for every prime is Brown \cite[\S2.1, Conjecture~1, p.~2; \S2.3, Theorem~2.1 and following paragraph, p.~4]{Brown2019}; the original theorem is \cite[Theorem~1.1 and its consequence, p.~950]{Brown2012}.

We spell out the completion step. In a separated complete positively filtered Lie algebra whose associated graded is free with only finitely many generators below each weight, choose lifts of the free graded generators. The resulting map from the weight-completed free Lie algebra induces an isomorphism on associated gradeds. Induction through successive finite weight quotients proves it is an isomorphism on each quotient; inverse limits give an isomorphism of the completed Lie algebras. A Levi splitting of the negatively weighted extension, as in \cite[Lemma~3.7, p.~9]{HainMatsumoto2001}, permits weight-homogeneous choices. Thus the completed universal enveloping algebra of $\operatorname{Lie}U$ is the weight-completed free associative algebra on the $s_m$. Its finite weight quotients are finite-dimensional.

Let $\ell_m:U(\mathbb Q_2)\to\mathbb Q_2(m)$ be the $m$th abelian coordinate, and put $a_m=\ell_m(\rho_U(h))$. Conjugation by $\gamma$ acts on that coordinate by $u^m$, because the unipotent part of a lift acts trivially on abelianization. Hence
\[
 \ell_m(\rho_U(h_x))=a_mu^{mx}.                             \tag{P2.5}
\]
Every $a_m$ is nonzero. Otherwise that coordinate would vanish on all the topological generators $h_x$, hence on $H$, contradicting the density in (P2.3) and the nonzero $m$th abelian quotient in (P2.4). No integral primitive normalization is asserted here.

\paragraph{Substitution at a finite weight cutoff.}
Suppose $g\neq1$ lies in the kernel of (P2.3). Let $r$ be its first nonzero Magnus degree, and let
\[
 0\neq\mu=\mu_r(g)\in\mathbb Z_2[[\mathbb Z_2^r]].
\]
In the completed free associative algebra put $Y_x=\rho_U(h_x)-1$. Its constant term is zero and its part of generator length one is
\[
 \sum_{m\geq3,\ m\text{ odd}}a_mu^{mx}s_m.
\]
Here generator length is distinct from arithmetic weight and from geometric depth.

Substitution $\delta_x\mapsto Y_x$ is well-defined and continuous coefficient by coefficient. At target weight at most $N$, only Magnus degrees at most $\lfloor N/3\rfloor$ contribute. The remaining coefficient functions on compact powers of $\mathbb Z_2$ are continuous and bounded in $\mathbb Q_2$. Integration against a bounded $\mathbb Z_2$-measure is defined by uniform approximation by locally constant functions, after scaling into $\mathbb Z_2$ if necessary. The same approximation proves continuity, multiplicativity and Fubini's identity. This construction needs no lattice uniform across all weights. The substitution and $\rho_U$ agree on every $h_x$ and are continuous homomorphisms on each finite weight quotient, so topological generation makes them agree on $H$. No universal property of a free pro-$2$ group is applied to a noncompact target.

The coefficient of the ordered word $s_{m_1}\cdots s_{m_r}$ in the image of $g-1$ is therefore
\[
 \left(\prod_{j=1}^ra_{m_j}\right)
 \int_{\mathbb Z_2^r}\prod_{j=1}^ru^{m_jx_j}\,d\mu(x).
\]
Lower Magnus degrees are zero, and higher ones have generator length greater than $r$. Since $\rho_U(g)=1$ and every $a_m\neq0$, we obtain
\[
 \int_{\mathbb Z_2^r}\prod_{j=1}^ru^{m_jx_j}\,d\mu(x)=0
 \quad(m_j\geq3\text{ odd}).                              \tag{P2.6}
\]
The variables remain ordered; this is not a symmetrized moment.

\paragraph{Analytic separation.}
The Amice transform
\[
 A_\mu(T_1,\ldots,T_r)
   =\int_{\mathbb Z_2^r}\prod_{j=1}^r(1+T_j)^{x_j}\,d\mu(x)
\]
lies in $\mathbb Z_2[[T_1,\ldots,T_r]]$ and converges on the open unit polydisc. It is injective as a transform of bounded measures: its coefficients are the moments against products of binomial polynomials, the multivariate Mahler basis. Equivalently it is the canonical isomorphism of the multivariable Iwasawa algebra with the power-series algebra.

For odd $m\geq3$ put $t_m=u^m-1$. These points lie in $4\mathbb Z_2$, and the distinct points $t_{1+2^j}$ converge to $u-1$. Distinctness follows from the infinite order of $u$. Equation (P2.6) says that $A_\mu$ vanishes on the full Cartesian grid of these points. The one-variable nonarchimedean identity theorem applies on the closed disc of radius $1/2$, contained in the radius-one convergence disc. Fixing grid values in all but one variable and applying that theorem successively in each variable shows $A_\mu=0$. Therefore $\mu=0$, a contradiction.

It follows that $\rho_U$ is injective on $H$. The kernel of the pro-$2$ tripod action is contained in $H$ because geometric abelianization records $\chi$. Its elements are also trivial under the induced rational outer action, so they are $1$. This proves Theorem~\ref{thm:pro2-kernel}.

\paragraph{What this proof does and does not decide.}
The proof uses the explicit arithmetic presentation $C_2*\mathbb Z_2$, not merely abstract freeness of $H$. It settles the field equality with the stated convention at infinity. It does not prove integral odd generation of the weight Lie algebra, and it does not transfer to irregular primes without a substitute for its linear-moment separation. The latter obstruction is made explicit below.


\subsection{Four meanings of an integral generator}
\label{sec:four-integral-meanings}

The following claims differ and must be kept separate:
\begin{enumerate}
\item[(A)] A nonzero integral odd class $\sigma_m\in\mathfrak g_{p,m}$ exists.
\item[(B)] The full-domain character
\[
 \kappa_m:G_{\mathbb Q(\mu_{p^\infty})}^{\mathrm{ab}}
       \longrightarrow\mathbb Z_p(m)
\]
is onto.
\item[(C)] An element of the weight subgroup $F^mG_p$ has $\kappa_m$-value $1$ in a fixed integral normalization.
\item[(D)] The map $\psi$ above is an isomorphism of graded $\mathbb Z_p$-Lie algebras.
\end{enumerate}

For odd $p$, (A) is already supplied by Soul\'e nonvanishing and Sharifi's integral construction \cite[Proposition~2.3, pp.~5--6]{Sharifi2002}. It is primitive relative to the \emph{actual image ideal}, which need not be $\mathbb Z_p$. At $p=2$, existence also follows from Brown and the actual lattice: multiply a nonzero rational odd generator by a sufficiently large power of $2$ to put it in $\mathfrak g_{2,m}$, then choose a representative in $F^mG_2$. This establishes existence, independently of the kernel theorem and of step~17.

For odd $p$, all of (B), for odd $m\geq3$, hold if and only if Vandiver holds at $p$. More precisely, each character is nonzero, and it is onto if $m\equiv1\pmod{p-1}$; simultaneous surjectivity is equivalent to Vandiver \cite[Theorem~1.1(1)--(3), p.~2]{IshiiKummer2024}. The theorem is attributed there to the classical Soul\'e and Ichimura--Sakaguchi work. We verified this precise primary-paper locator; we did not verify an original Ichimura--Sakaguchi theorem number and do not invent one. On all of $G_{\mathbb Q}$ the class is a cocycle, while on the cyclotomic kernel it is an ordinary continuous homomorphism.

There is a decisive restriction issue in (C). Sharifi's integral stable-derivation comparison gives, on $F^mG_p$,
\[
 \kappa_m=(p^{m-1}-1)(m-1)!\lambda_m^{(p)},
 \qquad \lambda_m^{(p)}\text{ integral}
 \quad\cite[(2.10), p.~7]{Sharifi2002}.
\]
Thus
\[
 \kappa_m(F^mG_p)\subseteq(m-1)!\mathbb Z_p(m).              \tag{I.1}
\]
At $p=3,m=5$, this image is contained in $3\mathbb Z_3$. Yet $3$ is regular, so (B) and (D) hold. Consequently full-domain surjectivity, weight-subgroup primitivity, and integral generation cannot be conflated. Formula (I.1) does not dispute the cohomological index formula for the full-domain cyclotomic class; it concerns a further restriction.

For (D), odd regular primes are settled positively by Sharifi--Brown. The prime $37$ gives a negative answer despite Vandiver and the corresponding vanishing of all the $K_{4k}$-groups at that prime. The proof follows.

\subsection{A counterexample to integral odd generation at \texorpdfstring{$37$}{37}}
\label{sec:integral37}

Let $\mathfrak h_p=\operatorname{im}\psi$. McCallum--Sharifi's notation in \S9 agrees with this choice: their ``Deligne conjecture in degree $i$'' means injectivity of $\operatorname{gr}^i\psi$ \cite[p.~24]{McCallumSharifi2003}. Brown gives precisely this injectivity, because the rational map is injective and every graded piece of the integral free source is torsion-free.

\begin{proposition}[T-INTEGRAL37]\label{prop:37-counterexample}
At $p=37$, Vandiver holds, all $K_{4k}(\mathbb Z)\otimes\mathbb Z_{37}$ vanish for $k\geq1$, and all full-domain odd Soul\'e characters are onto. Nevertheless
\[
 (\mathfrak g_{37}^{\mathrm{ab}})_{32}\neq0,
 \qquad
 \mathfrak h_{37,32}\subsetneq\mathfrak g_{37,32}.
\]
In particular $\mathfrak g_{37}$ cannot be generated over $\mathbb Z_{37}$ by homogeneous elements all of odd weight, and is not the free graded integral Lie algebra on one generator in every odd weight at least three.
\end{proposition}

\begin{proof}
The only irregular index for $37$ is $32$. Vandiver at $37$ can be verified by the small exact modular certificate accompanying this report, using the criterion of Buhler--Harvey \cite[\S4, pp.~4--5]{BuhlerHarvey2011}: take $q=149$ and $z=16$, with $z$ of exact order $37$. The required product is $V=48\pmod{149}$ and $V^{(149-1)/37}=142\not\equiv1\pmod{149}$. The computation uses exact integer arithmetic and was independently cross-checked; it is not a formal proof certificate. The modern K-theoretic comparison recorded elsewhere in this report gives the asserted $K_{4k}$ vanishing. The full-domain character assertion follows from \cite[Theorem~1.1(3)]{IshiiKummer2024}.

McCallum--Sharifi's Theorem~7.5 proves that their pairing at index $32$ for $37$ is nonzero. Their Corollary~8.3, p.~23, makes a coefficient explicit: $e_{3,32}$ is nonzero, represented by $-11$ up to a nonzero normalization scalar. Their Theorem~9.7, p.~28, applies when Vandiver holds, $m$ is the least irregular even index, rational Deligne holds through degree $m$, and some $e_{i,m}\neq0$ has $i<m/2$. All hypotheses hold with $(p,m,i)=(37,32,3)$, by Brown and $3<16$.

For clarity, the proof of that cited theorem identifies the mechanism. In degrees below $32$, the integral map is bijective. Thus the weight-$32$ commutator part of $\mathfrak g$ equals that of $\mathfrak h$. The nonzero cup coefficient, through their Proposition~9.6, prevents bijectivity at $32$; Brown's injectivity makes the defect a failure of surjectivity. Their theorem consequently gives both displayed conclusions.

If a graded Lie algebra is generated by homogeneous odd-weight elements, its abelianization is generated by their odd-weight images and has no even-weight component. The nonzero degree-$32$ abelianization contradicts this. This proves the proposition.
\end{proof}

The authors' Magma computation in Theorem~7.5 is used as a published theorem; we did not rerun it. The deduction above, including every hypothesis of Theorem~9.7, was read in full. It does not say that $\mathfrak h_{37}$ is nonfree: rational injectivity makes it a free integral subalgebra. It says that its degree-$32$ lattice is proper. The graded pieces of $\mathfrak g_{37}$ remain torsion-free; their Lie abelianization can have $37$-torsion.

Most importantly, Proposition~\ref{prop:37-counterexample} does not refute $\Omega_{37}^*=\Omega_{37}$. An arithmetic source and its faithful image can have these same extra integral directions. This is a counterexample to the proposed integral-generation implication, not a counterexample to the kernel-field equality.

Sharifi also identifies the first possible defect under Vandiver: if $\Psi$ is not onto and $m$ is its first defective weight, then $m$ must be even and $p\mid B_m$ \cite[Theorem~4.1, p.~10]{Sharifi2002}. The odd-weight alternative would violate Vandiver. This explains why removing the plus class-group obstruction does not eliminate all integral-generation obstructions.

\subsection{Irregular primes: the sharper remaining gate}
\label{sec:irregular-gate}

Write $H_p=\operatorname{Gal}(\Omega_p/K)$ and $M_p=H_p^{\mathrm{ab}}$, with the conjugation action of $\Gamma=\operatorname{Gal}(K/F)\simeq\mathbb Z_p$ and $\Lambda=\mathbb Z_p[[\Gamma]]$. A useful correction is that $H_p$ is already free pro-$p$ even at irregular primes. Iwasawa's original Theorem~2 states that the maximal $S$-ramified pro-$p$ group over a cyclotomic $\mathbb Z_p$-field is free exactly when its $\mu$-invariant is zero \cite[pp.~273--274]{Iwasawa1981}. Ferrero--Washington applies to $K/\mathbb Q$, which is abelian. This does not make the finite-base group $\operatorname{Gal}(\Omega_p/F)$ free.

Let $A_\infty$ be the direct limit, under extension of ideals, of the $p$-parts of the class groups in the cyclotomic tower. Kummer duality gives
\[
 M_p^+\simeq(A_\infty^-)^\vee(1)
 \quad\cite[Corollary~3.3.14 and Proposition~3.4.7]{SharifiIwasawaNotes}.
                                                               \tag{I.2}
\]
The finite-level minus class groups inject into $A_\infty^-$, by the same source's Proposition~3.4.1. At an irregular prime that minus class group is nonzero, so $M_p^+\neq0$. With $\mu=0$, this is a nonzero finite-rank free $\mathbb Z_p$-module but a torsion $\Lambda$-module. The distinction between $p$-torsion and $\Lambda$-torsion is essential.

Every odd linear Soul\'e coordinate kills $M_p^+$: complex conjugation acts by $+1$ on the source and $-1$ on the target $\mathbb Q_p(m)$, so equivariance gives twice the value equal to zero. Thus the degree-one separation used in the pro-$2$ proof fails on a concrete nonzero module. A lift of a nonzero class in $M_p^+$ can have a nonzero first source Magnus coefficient while all target generator-length-one coordinates vanish. Its image may still be nontrivial at higher length. Abstract freeness of $H_p$, weak Leopoldt, or $\mu=0$ alone therefore cannot repeat that proof.

A sharper continuation is to control the full nonlinear lifts of the $\Gamma$-relations defining this torsion module. Their first nonlinear terms are cup products; the degree-$32$ result at $37$ detects one such direction through a bracket. Kernel equality requires detection of every remaining direction and its iterated commutators, or an equivalent separation theorem for the full noncommutative relation ideal. A class-group order or characteristic ideal alone does not specify those nonlinear terms. The following subsections supply an all-weight separation proof at $37$ and its generalization under Vandiver and odd cup-product surjectivity. In particular the field equality holds for every odd prime below $1000$; the remaining all-prime arithmetic conditions are explicit. Neither Proposition~\ref{prop:37-counterexample} nor (I.2) gives a negative answer.

% Checked deduction: author, root, vandiver and archive whole-proof readings completed.
% Candidate and first independent audit: work/irregular37_pole_candidate.md.
\subsection{A nonlinear separation argument at \texorpdfstring{$37$}{37}}
\label{sec:kernel37-poles}

The failure of integral odd generation at $37$ does not decide the kernel-field question. The following argument detects its extra even arithmetic direction through a pole in a quadratic moment function. It uses the published cup-product computation, but supplies a separate group-theoretic and analytic proof of faithfulness. The proof is a deduction developed in this run, not a theorem located in the literature or a consequence of step~17.

\begin{theorem}[K17-37-KERNEL]\label{thm:kernel37-poles}
For the maximal pro-$37$ extension $\Omega_{37}$ of $\mathbb Q(\mu_{37})$ unramified outside $37$, the tripod kernel field satisfies
\[
 \Omega_{37}^*=\Omega_{37}.
\]
This is compatible with $(\mathfrak g_{37}^{\mathrm{ab}})_{32}\ne0$ and failure of integral odd generation.
\end{theorem}

We first prove the separation statement needed for the arithmetic application.

\paragraph{K17-POLE-ONE: the compact orbit basis and separation criterion.}
Let $p$ be odd, let $\Gamma=\langle\gamma\rangle\simeq\mathbb Z_p$, and put $\Lambda=\mathbb Z_p[[\Gamma]]$. Suppose $H$ is free pro-$p$ with a continuous $\Gamma$-action, and that its abelianization has the compatible decomposition
\[
 H^{\mathrm{ab}}=
 \left(\bigoplus_{a\in A}\Lambda e_a\right)\oplus\mathbb Z_p e_y,
 \qquad \gamma e_y=\alpha e_y,\qquad \alpha\in1+p\mathbb Z_p.       \tag{N.1}
\]
Here $A$ consists of all odd residue classes modulo $p-1$. A commuting finite group $\Delta$ of order $p-1$ acts on $e_a$ by character $a$ and on $e_y$ by an even character. Choose lifts $x_a,y\in H$, and set
\[
 x_{a,t}=\gamma^t x_a\gamma^{-t},\qquad t\in\mathbb Z_p.
\]
These are a free compact generating space for $H$, in the following precise sense.

Let $F_X$ be the inverse limit of the finite-rank free pro-$p$ groups on
\[
 X_n=(A\times\mathbb Z/p^n)\amalg\{y\},
\]
with maps folding the residue indices. This limit has the universal property for continuous maps from $X=(A\times\mathbb Z_p)\amalg\{y\}$ into a pro-$p$ group: after composition with any finite quotient of the target, such a map factors through an $X_n$. Restricting that property to abelian targets gives
\[
 F_X^{\mathrm{ab}}=\mathbb Z_p[[X]]=
 \Lambda^{|A|}\oplus\mathbb Z_p.
\]
Thus the natural map $F_X\to H$ is an isomorphism on abelianizations. It is onto: a proper closed subgroup of a pro-$p$ group lies in a proper open subgroup and then in a maximal open subgroup, so cannot surject on the Frattini quotient. Since $H$ is free pro-$p$, projectivity supplies a section. Its closed image surjects on $F_X/\Phi(F_X)$ and must therefore be all of $F_X$. The map is an isomorphism. This argument uses neither infinite-rank Hopficity nor an assumed interchange of abelianization and inverse limits.

The compatible finite-rank Magnus embeddings consequently give an injective measure-valued Magnus expansion of $H$. Assign degree $1$ to the orbit letters and degree $2$ to $y$. Every nontrivial element has a first finite nonzero degree. In any degree, each word pattern has a bounded $\mathbb Z_p$-measure as coefficient in the orbit parameters of its letters of degree $1$.

\paragraph{A quadratic meromorphic moment.}
Suppose there is a continuous dense map $\rho:H\to U(\mathbb Q_p)$ into the actual rational tripod image, whose Lie algebra is the completed free Lie algebra on odd generators $s_m$, $m\ge3$. Suppose $\gamma$ is represented by a Levi element with cyclotomic value $u$, a topological generator of $1+p\mathbb Z_p$. Thus it scales $s_m$ by $u^m$, and $\Delta$ acts by character $m$.

Let $\ell_m$ be the $m$th abelian coordinate. Equivariance gives
\[
 \ell_m(\rho(x_{a,t}))=
 \begin{cases}a_m u^{mt},&m\equiv a\pmod{p-1},\\0,&\text{otherwise}.
 \end{cases}
\]
Each $a_m$ is nonzero by density and the compact generating space: $y$ has no odd linear coordinate by parity. Rescale each free letter by $a_m$, weight by weight over $\mathbb Q_p$. The linear part of $\rho(x_{a,t})-1$ now has coefficient $u^{mt}$ on its matching letters. No uniformly integral change of lattice is asserted.

Set
\[
 E=\gamma y\gamma^{-1}y^{-\alpha}\in\overline{[H,H]}.
\]
Its degree-two coefficients involving two orbit letters give bounded measures $\mu_{a,b}$ and transforms
\[
 C_{a,b}(T_1,T_2)=
 \int_{\mathbb Z_p^2}(1+T_1)^t(1+T_2)^v\,d\mu_{a,b}(t,v)
 \in\mathbb Z_p[[T_1,T_2]].
\]
Assume that for one ordered pair $(a_0,b_0)$,
\[
 C_{a_0,b_0}(0,0)\in\mathbb Z_p^\times.                         \tag{N.2}
\]
We will verify this hypothesis arithmetically at $37$.

Let $B_{m,n}$ be the coefficient of $s_ms_n$ in $\rho(y)-1$. Since the latter has no length-one part, comparison in length two yields
\[
 (u^{m+n}-\alpha)B_{m,n}
   =C_{a,b}(u^m-1,u^n-1),\qquad a\equiv m,\ b\equiv n\pmod{p-1}.
                                                                    \tag{N.3}
\]
Indeed the two sides of $E=\gamma y\gamma^{-1}y^{-\alpha}$ give the factor $u^{m+n}-\alpha$, while only the two-orbit-letter Magnus terms can contribute to target length two. Terms of $\overline{[H,H]}$ involving $y$ have assigned degree at least three.

Write
\[
 D(T_1,T_2)=(1+T_1)(1+T_2)-\alpha.
\]
Equation (N.3) represents the length-two kernel of $y$ by $C_{a,b}/D$. Divisions will only be performed in a fraction field after clearing denominators. The numerator in (N.2) is a unit. The divisor is genuine in the open unit polydisc, since $T_2=\alpha/(1+T_1)-1$ and $\alpha-1\in p\mathbb Z_p$.

\paragraph{Ordered tilings and pole separation.}
Suppose $g\ne1$ belongs to $\ker\rho$, and let $d$ be its first nonzero assigned Magnus degree. A source word of degree $d$ specifies a subset
\[
 J\subseteq\{1,\ldots,d-1\},\qquad
 \text{no two elements of }J\text{ adjacent}.
\]
Each $j\in J$ denotes one $y$ filling target positions $j,j+1$; the other positions are singleton orbit letters. For each pattern and its singleton types, let $A_J$ be the Amice transform of its coefficient measure. It depends only on the unpaired variables.

Fix the residue classes of an ordered target word of length $d$. Its vanishing coefficient gives, at $T_i=u^{m_i}-1$,
\[
 \sum_J A_J(T_i:i\text{ unpaired})
       \prod_{j\in J}
       \frac{C_{a_j,a_{j+1}}(T_j,T_{j+1})}{D_j}=0,
 \qquad D_j=(1+T_j)(1+T_{j+1})-\alpha.                         \tag{N.4}
\]
There are finitely many tilings. Multiply by all $D_j$, $1\le j<d$. The resulting integral analytic identity holds even at grid points where a denominator vanishes, by (N.3). Smaller source degrees vanish by choice of $d$, and larger ones have target length greater than $d$; thus no omitted source terms contribute.

The finite-weight substitution is valid exactly as in the pro-$2$ proof. At every arithmetic weight cutoff only finitely many source degrees contribute. The coefficient functions on compact orbit spaces are bounded and continuous; integration, uniform locally constant approximation, and Fubini give a continuous multiplicative map. It agrees with $\rho$ on the generators and hence on $H$.

If $\alpha=u^w$ for a positive integer $w$, restrict all test weights to exceed $w$; otherwise no exclusion is needed. This removes every possible zero denominator on the test grid. For each odd residue class $a$, the remaining tails $m=a+(p-1)p^n$ still give distinct points converging to $u^a-1$ in $p\mathbb Z_p$. Applying the analytic identity theorem successively in all variables makes the denominator-cleared version of (N.4) an identity of power series. Hence (N.4) holds in the fraction field of
\[
 R=\mathbb Z_p[[T_1,\ldots,T_d]][1/p].
\]
This ring has a uniform bound on coefficient denominators. Replacing it by $\mathbb Q_p[[T_1,\ldots,T_d]]$ would destroy the argument by allowing inverses of $D_j$ with unbounded denominators.

Choose a nonzero coefficient pattern with the maximum number of $y$ letters, and denote its tiling by $J$. Assign target types $(a_0,b_0)$ at each of its pairs and the specified source type at each singleton. This fixes an identity (N.4) with $A_J\ne0$. Consider
\[
 I_J=(D_j:j\in J)\subset R.
\]
This is prime: its disjoint equations substitute $T_{j+1}=\alpha/(1+T_j)-1$, leaving a power-series ring in the unpaired variables and one variable from each pair, with $p$ inverted. The substitutions are legitimate in the $(p,T)$ topology because their constant term is in $p\mathbb Z_p$.

For $k\notin J$, $D_k\notin I_J$. Its neighboring positions belong to different blocks; the remaining variables of these blocks are independent, and their product is not identically $\alpha$. Thus $D_k$ is a unit in $R_{I_J}$. Multiply (N.4) by $\prod_{j\in J}D_j$ and reduce modulo $I_J$ in this localization. A tiling omitting any pair of $J$ vanishes. A tiling properly containing $J$ has more $y$ letters and its coefficient is zero by maximality. Only $J$ remains:
\[
 A_J\prod_{j\in J}C_{a_0,b_0}(T_j,T_{j+1})=0
 \quad\text{in }\operatorname{Frac}(R/I_J).
\]
Every $C$ factor is a unit. The unpaired-variable ring embeds in $R/I_J$, so $A_J$ remains nonzero, a contradiction. If $J$ is empty, maximality leaves just its singleton coefficient, giving the same contradiction directly. This proves injectivity under (N.1)--(N.2).

\paragraph{The arithmetic module at $37$.}
Now set $p=37$, $F=\mathbb Q(\mu_{37})$, $K=\mathbb Q(\mu_{37^\infty})$, $P=\operatorname{Gal}(\Omega_{37}/F)$, and $H=\operatorname{Gal}(\Omega_{37}/K)$. Choose commuting lifts of $\Delta=\operatorname{Gal}(F/\mathbb Q)$ and $\Gamma=\operatorname{Gal}(K/F)$ using Sharifi's procyclic construction \cite[\S2, p.~3]{Sharifi2002}. If the chosen generator has cyclotomic value $u_0$, replace it by its $v$th power, where $v=\log(1+37)/\log(u_0)\in\mathbb Z_{37}^{\times}$. The resulting $\gamma$ still commutes with $\Delta$ and has $\chi(\gamma)=1+37$.

Iwasawa's theorem and Ferrero--Washington give freeness of $H$ \cite[Theorem~2, pp.~273--274]{Iwasawa1981}. Put $M=H^{\mathrm{ab}}$. Weak Leopoldt gives $\operatorname{rank}_{\Lambda}M=18$. Kummer duality gives that $M^+$ is $\Lambda$-torsion \cite[Theorems~3.3.3--3.3.4, Proposition~3.4.7]{SharifiIwasawaNotes}. Also $M$ is $\mathbb Z_{37}$-torsion-free because $H$ is free pro-$37$.

The minimal generators in McCallum--Sharifi, \S8, pp.~22--23, consist of $\gamma$, eighteen odd $x_m$ for $3\le m\le37$, and a single even generator $y$ of character $32$ \cite{McCallumSharifi2003}. Vandiver at $37$ and its unique irregular index $32$ were checked above. The exact sequence
\[
 0\longrightarrow M_{\Gamma}\longrightarrow P^{\mathrm{ab}}
   \longrightarrow\Gamma\longrightarrow0
\]
follows from the split semidirect product $P=H\rtimes\Gamma$. Consequently $M/(37,T)M$ has exactly eighteen odd directions and one even direction. Nakayama gives eighteen generators for $M^-$. Its rank is eighteen, since $M^+$ has rank zero. A surjection $\Lambda^{18}\to M^-$ has rank-zero kernel, which must vanish as a submodule of a free module over the domain $\Lambda$. Thus
\[
 M^-\simeq\bigoplus_{a\text{ odd}\bmod36}\Lambda e_a.           \tag{N.5}
\]

Similarly $M^+$ is cyclic, generated by the class of $y$, and is nonzero because its quotient modulo $(37,T)$ has dimension one. The explicit relation in McCallum--Sharifi, Corollary~8.3, p.~23, begins
\[
 37y-3[\gamma,y]-11[x_3,x_{29}]-[x_7,x_{25}]+\cdots=0
                                                                  \tag{N.6}
\]
in descending $p$-central degree two. Linearizing the full relator in the cyclotomic kernel and projecting onto the plus part gives an annihilator
\[
 f(T)y=0,\qquad f(T)\equiv37-3T\pmod{(37,T)^2}.                \tag{N.7}
\]
Here is the filtration justification for the coefficient congruence. Before imposing arithmetic relations, linearize the finite-rank free presentation group into $A_{\mathrm{free}}\rtimes\Gamma$, where $A_{\mathrm{free}}$ is the free $\Lambda$-module on the non-cyclotomic generators. Its filtration
\[
 (p,T)^{n-1}A_{\mathrm{free}}\rtimes p^{n-1}\Gamma
\]
respects commutators and $p$th powers; use
$(1+T)^{p^{n-1}}-1\in(p,T)^n$. Thus descending $p$-central degree three contributes only $(p,T)^2$ to each coefficient. Read the $y$ coefficient there, then pass to $M^+$. The odd commutators vanish and odd linear terms disappear in this projection.

Weierstrass preparation makes $f$ a unit times $T-\beta$, where $\beta\in37\mathbb Z_{37}$ and $v_{37}(\beta)=1$. Hence $M^+$ is a quotient of $\Lambda/(f)\simeq\mathbb Z_{37}$. It is nonzero and $\mathbb Z_{37}$-torsion-free, so
\[
 M^+\simeq\mathbb Z_{37},\qquad\gamma e_y=(1+\beta)e_y.        \tag{N.8}
\]
Equations (N.5)--(N.8) prove (N.1), with $\alpha=1+\beta$.

The lifts can preserve the preceding Frattini classes: project their abelian classes with the integral $\Delta$-idempotents and lift back through $H\to M$. The projections fix their prescribed residual characters. Each odd projected class generates its free rank-one $\Lambda$-eigenspace, and the even class generates $M^+$. A change of finite free presentation generators with the same Frattini classes is induced by an automorphism: the endomorphism is surjective by the Frattini criterion and injective by Hopficity of a finitely generated profinite group. Such a change fixes descending $p$-central degree two: both $p$th powers and commutators retain their classes. Unit rescalings within the one-dimensional eigenspaces only rescale the displayed nonzero coefficients by units. Thus these changes preserve the information used from (N.6).

\paragraph{Recovering the unit quadratic numerator.}
For the compact basis supplied by (N.1), write $\varphi$ for the $\Gamma$-automorphism of $H$ and $E=\varphi(y)y^{-\alpha}$. The group $P$ has the presentation on $\gamma,x_a,y$ with one defining relation
\[
 R=\gamma y\gamma^{-1}\varphi(y)^{-1}.
\]
To verify the presentation, assign images to these generators in any pro-$p$ target satisfying $R$. Send the compact orbit generator $x_{a,t}$ to $\gamma^t x_a\gamma^{-t}$. The compact free universal property gives a map from $H$; its $\Gamma$-compatibility is automatic on the orbits and is exactly $R$ on $y$. Hence it extends to $H\rtimes\Gamma$.

In descending $p$-central degree two the relation has class
\[
 [\gamma,y]+(1-\alpha)y-E_2,
\]
where $E_2$ is the quadratic Magnus term after folding the orbit parameters to their augmentations. Its ordered $x_a x_b$ coefficient is, up to the commutator sign convention, $C_{a,b}(0,0)$. Because there is one defining relation, any relation's degree-two initial term is a scalar multiple of this one: conjugation fixes that class and $p$th powers of the relator vanish in that degree. No assertion about freeness of the full relation module is required.

In (N.6) the coefficients of $[\gamma,y]$ and $[x_3,x_{29}]$ are both units modulo $37$. It follows that $C_{3,29}(0,0)$ is a unit. This is (N.2), with $(a_0,b_0)=(3,29)$.

\paragraph{The actual rational image and conclusion.}
Hain--Matsumoto's actual-image comparison and Brown's rational Deligne--Ihara theorem supply the dense map $H\to U(\mathbb Q_{37})$ and its completed free odd Lie algebra \cite[Proposition~8.3, Theorem~8.4]{HainMatsumoto2001}\cite[\S2.1--2.3]{Brown2019}. Arrange that the image of $\gamma$ lies in a Levi torus by successively solving in positive weights, since $u^w-1\ne0$ in $\mathbb Q_{37}$ for every $w>0$. This converges in the weight completion. The commuting finite group $\Delta$ then also lies in that torus: its possible unipotent part would centralize $\gamma$, whose unipotent centralizer is trivial by the same weight argument.

All hypotheses of the separation argument are now verified, so the actual rational action is injective on $H$. The kernel of the pro-$37$ tripod action lies in $H$ because abelianization records the cyclotomic character, and it is trivial under the induced rational action. Thus it is trivial on $P$, proving Theorem~\ref{thm:kernel37-poles}.

This proof deals with the single even arithmetic direction at $37$. It does not supply an all-irregular-primes theorem. With several even directions or higher $\Lambda$-torsion rank, the corresponding matrix and higher-order pole separation must be proved separately.

% CHECKED DEDUCTION: full self-reading, two independent whole-proof audits, and parent review passed.
% Existing checked p=2 and p=37 arguments remain separate.
\subsection{Several even directions and the range \texorpdfstring{$p<1000$}{p<1000}}
\label{sec:multieven-poles}

The preceding pole argument extends to several even arithmetic directions and to characteristic polynomials with repeated roots. The arithmetic hypothesis is most cleanly stated directly in Galois cohomology.

\begin{theorem}[K17-POLE-GENERAL]\label{thm:pole-general}
Let $p$ be an odd prime satisfying Vandiver, let $F=\mathbb Q(\mu_p)$, and let $P$ be the Galois group of its maximal pro-$p$ extension unramified outside $p$. Assume that the cup products of the odd-character summand of $H^1(P,\mathbb F_p)$ span $H^2(P,\mathbb F_p)$. Then
\[
 \Omega_p^*=\Omega_p.
\]
No assumption that the relevant Iwasawa polynomials have degree one, split over $\mathbb Q_p$, or have distinct roots is required.
\end{theorem}

Here odd means the $-1$ eigenspace for complex conjugation, equivalently the sum of the odd Teichm\"uller characters. Cup/relator duality says that the hypothesis supplies a nonzero odd--odd commutator coefficient in each nonzero even relation eigenspace. This is a statement about the given arithmetic group and its full cup product; it does not require a chosen basis of cyclotomic units.

\paragraph{The arithmetic module and its companion bases.}
Set $K=\mathbb Q(\mu_{p^\infty})$, $H=\operatorname{Gal}(\Omega_p/K)$, $\Gamma=\operatorname{Gal}(K/F)=\langle\gamma\rangle$, $\Lambda=\mathbb Z_p[[\Gamma]]$, and $M=H^{\mathrm{ab}}$. Choose $\gamma$ commuting with a lift of $\Delta=\operatorname{Gal}(F/\mathbb Q)$ and with $\chi(\gamma)=u=1+p$, as above.

Iwasawa's theorem and Ferrero--Washington make $H$ free pro-$p$, so $M$ is $\mathbb Z_p$-torsion-free. Weak Leopoldt gives $\operatorname{rank}_{\Lambda}M=(p-1)/2$. Kummer duality makes $M^+$ a torsion $\Lambda$-module; it is finitely generated over $\mathbb Z_p$ because it is pseudo-isomorphic to $(X_\infty^-)^\iota(1)$, whose $\mu$-invariant is zero. Finite kernels and cokernels preserve this finite generation. Thus $M^+$ is finite free over $\mathbb Z_p$ \cite[Theorem~2]{Iwasawa1981}\cite[Theorems~3.3.3--3.3.4, Proposition~3.4.7, Corollary~3.4.8]{SharifiIwasawaNotes}.

Under Vandiver the minimal generators of $P$ in McCallum--Sharifi, \S8, pp.~22--23, are $\gamma$, one $x_a$ for each odd residue class $a$, and one $y_e$ for each occurring even character $e$ \cite{McCallumSharifi2003}. From
\[
 0\longrightarrow M_\Gamma\longrightarrow P^{\mathrm{ab}}
  \longrightarrow\Gamma\longrightarrow0
\]
and Nakayama, $M^-$ has $(p-1)/2$ generators over $\Lambda$ and the same rank. The kernel of the resulting surjection from a free module has rank zero and is zero. Hence
\[
 M^-\simeq\bigoplus_{a\text{ odd}\bmod p-1}\Lambda e_a.       \tag{G.1}
\]
Each nonzero even component $T_e=M^{(e)}$ is cyclic over $\Lambda$ and finite free over $\mathbb Z_p$. Put $\lambda_e=\operatorname{rank}_{\mathbb Z_p}T_e$ and choose a cyclic generator $v_e$. If $q_e(z)$ is the characteristic polynomial of $\gamma$ on $T_e$, then
\[
 q_e(z)\equiv(z-1)^{\lambda_e}\pmod p.                         \tag{G.2}
\]
Indeed the continuous action of the pro-$p$ group $\Gamma$ on $T_e/pT_e$ has finite $p$-group image and is unipotent.

The vectors $v_e,\gamma v_e,\ldots,\gamma^{\lambda_e-1}v_e$ form a $\mathbb Z_p$-basis. Cayley--Hamilton makes their span stable under $\gamma$ and $\gamma^{-1}$, since $q_e$ has unit constant coefficient. The span is compact and closed, so contains the entire $\Lambda$-span of $v_e$. It therefore surjects onto $T_e$, and a surjection between free $\mathbb Z_p$-modules of the same finite rank is an isomorphism.

Choose actual lifts $x_a,y_e\in H$ and set
\[
 x_{a,t}=\gamma^t x_a\gamma^{-t},\quad t\in\mathbb Z_p,
 \qquad y_{e,j}=\gamma^j y_e\gamma^{-j},\quad0\le j<\lambda_e.
\]
These constitute a free compact basis of $H$. To recall the complete argument, take the inverse limit of the free pro-$p$ groups on the finite quotients of this compact generating space. Its abelianization is precisely (G.1) together with the indicated finite bases of the $T_e$. The natural map to $H$ is an isomorphism on abelianizations, hence is onto by the Frattini criterion for closed subgroups. Projectivity of $H$ supplies a section; its closed image is Frattini-surjective in the source and hence equals the source. The map is an isomorphism.

We may make the lifts exact nonlinear $\Delta$-eigenlifts by Sharifi's averaging lemma \cite[Lemma~2.1, pp.~3--4]{Sharifi2002}:
\[
 \delta x_a\delta^{-1}=x_a^{\omega(\delta)^a},\qquad
 \delta y_e\delta^{-1}=y_e^{\omega(\delta)^e}.
\]
This preserves the prescribed abelian eigenspace classes. Commutation of $\gamma$ and $\Delta$ gives the same eigenproperty for all their indicated conjugates.

\paragraph{The companion error and its cup coefficient.}
Write
\[
 q_e(z)=z^{\lambda_e}-\sum_{j=0}^{\lambda_e-1}c_{e,j}z^j
\]
and, with the product taken in increasing $j$, define
\[
 E_e=(\gamma^{\lambda_e}y_e\gamma^{-\lambda_e})
          \left(\prod_jy_{e,j}^{c_{e,j}}\right)^{-1}.
\]
The characteristic-polynomial relation in $M$ puts $E_e$ in the closed commutator subgroup of $H$. Its quadratic Magnus coefficients in two odd orbit letters give bounded measures and integral Amice transforms $C_{e;a,b}(T_1,T_2)$. We claim that for each $e$, there is an ordered odd pair $(a_e,b_e)$ such that
\[
 C_{e;a_e,b_e}(0,0)\in\mathbb Z_p^\times.                       \tag{G.3}
\]
The following relation-space argument supplies this comparison even when there are several even characters.

Let $F_{\mathrm{free}}$ be free pro-$p$ on $\gamma,x_a,y_e$, retaining only $j=0$ as an even generator. Substitute the actual conjugate words for all $y_{e,j}$ and $x_{a,t}$ in the compact free expansion of $E_e$. The group $P=H\rtimes\Gamma$ has one defining relator per even character:
\[
 R_e=\gamma^{\lambda_e}y_e\gamma^{-\lambda_e}
       \left(\prod_jy_{e,j}^{c_{e,j}}\right)^{-1}E_e^{-1}.       \tag{G.4}
\]
This is verified by the universal property. An assignment of $\gamma,x_a,y_e$ into a pro-$p$ target satisfying (G.4) gives images of all compact generators of $H$. Compatibility with $\gamma$ is automatic on the odd orbits and along each finite even chain except at its last element; the last compatibility is exactly (G.4). It therefore extends uniquely to the semidirect product.

Let $R$ be the closed relation subgroup. The exact eigenlifts let $\Delta$ act on $F_{\mathrm{free}}$ by fixing $\gamma$ and raising each other generator to its prescribed Teichm\"uller unit power. These are finite-order automorphisms; the quotient to $P$ is equivariant. Abelianize the kernel of $F_{\mathrm{free}}\to\Gamma$. Its free relative abelianization is
\[
 A_{\mathrm{free}}=
   \bigoplus_a\Lambda X_a\oplus\bigoplus_e\Lambda Y_e.
\]
The image of $R_e$ is exactly $q_e(\gamma)Y_e$. The word $E_e$ is a closed commutator word in the compact free basis, so has zero relative abelianization before imposing any arithmetic relations. Normal conjugation acts on this module through $\Gamma$. Hence the whole closed image is
\[
 R_{\mathrm{lin}}=\bigoplus_eq_e(\gamma)\Lambda Y_e,             \tag{G.5}
\]
a free $\Lambda$-module, with its indicated character $e$ on each summand.

Linearization gives a surjection
\[
 R/\overline{R^p[R,F_{\mathrm{free}}]}
       \longrightarrow R_{\mathrm{lin}}/(p,T)R_{\mathrm{lin}}.  \tag{G.6}
\]
The source is generated by the number of defining relators and the target has exactly that dimension. Thus (G.6) is an equivariant isomorphism, and the class of $R_e$ has character $e$. The cup-normalized character-$e$ relation is therefore a nonzero scalar multiple of $R_e$ in this quotient. All relators lie in the Frattini subgroup; passing to descending $p$-central degree two annihilates $R^p$ and $[R,F_{\mathrm{free}}]$.

The companion product in (G.4) contains only $y_e$ and its $\gamma$-conjugates and contributes no odd--odd bracket in that degree. Its odd--odd coefficient is consequently minus the augmentation coefficient of $E_e$, namely $-C_{e;a,b}(0,0)$ up to the commutator convention. The assumed odd cup surjectivity and the usual minimal-presentation cup/relator duality prove (G.3) \cite[\S8, pp.~21--23]{McCallumSharifi2003}. No freeness or asphericity claim about the full relation module is being used.

\paragraph{Quadratic target coefficients and partial fractions.}
Use the actual continuous dense rational representation $H\to U(\mathbb Q_p)$ and the completed free odd Lie algebra supplied by Hain--Matsumoto and Brown \cite{HainMatsumoto2001,Brown2019}. Conjugate rationally so that $\gamma$ is in a Levi torus; the successive positive-weight equations are solvable because $u^w-1\ne0$. The commuting $\Delta$ then also acts by its weight characters. Normalize each odd free letter $s_m$ so that the linear coefficient of $\rho(x_{a,t})-1$ on its matching residue class is $u^{mt}$. The necessary scalars are nonzero by density. Every $y_{e,j}$ has target generator length at least two by parity.

For an ordered pair $(m,n)$ put $z=u^{m+n}$, and let $B_{e,j;m,n}$ be the coefficient of $s_ms_n$ in $\rho(y_{e,j})-1$. The actual conjugate definition gives
\[
 B_{e,j;m,n}=z^jB_{e,0;m,n},
 \qquad
 q_e(z)B_{e,0;m,n}=C_{e;a,b}(u^m-1,u^n-1),                    \tag{G.7}
\]
where $a\equiv m$, $b\equiv n\pmod{p-1}$. The second equality is the length-two component of $E_e$. These coefficients have character $e$, so vanish when $a+b\not\equiv e$. For those residue classes (G.7) makes $C_{e;a,b}$ vanish on the entire Cartesian weight grid; the successive analytic identity theorem forces $C_{e;a,b}=0$ as a power series. Thus wrong-character numerators vanish identically, and the unit pair in (G.3) has $a_e+b_e\equiv e$.

Before dividing, exclude all exceptional weight sums. There are finitely many positive integers $w$ such that $u^w$ is a root of some $q_e$, since $u$ has infinite order. Restrict every test weight $m$ to exceed all such $w$. Then all pair denominators are nonzero. In each fixed odd residue class, the tails $m=a+(p-1)p^n$ still accumulate $p$-adically at $a$, and suffice for every analytic identity below.

Let $L/\mathbb Q_p$ be a finite extension splitting all $q_e$, with valuation ring $\mathcal O_L$ and uniformizer $\pi$. Every root $\alpha$ satisfies $\alpha-1\in\pi\mathcal O_L$ by (G.2). If its multiplicity is $n_{e,\alpha}$, ordinary partial fractions give the two bases of the same space of strictly proper rational functions:
\[
 \left\{\frac{z^j}{q_e(z)}:0\le j<\lambda_e\right\},
 \qquad
 \left\{\frac1{(z-\alpha)^k}:1\le k\le n_{e,\alpha}\right\}_{\alpha}.
\]
Apply the resulting constant invertible $L$-linear change to the degree-two Magnus letters $y_{e,j}$. This is a linear change in the associative coefficient algebra, not a change of actual pro-$p$ generators. At each fixed degree, only finitely many scalar coefficients occur, so measures retain a finite common denominator bound. The first nonzero assigned degree is preserved.

In the new letters $Y_{e,\alpha,k}$, the leading target kernels are
\[
 \frac{C_{e;a,b}(T_1,T_2)}
      {\bigl((1+T_1)(1+T_2)-\alpha\bigr)^k}.                   \tag{G.8}
\]
All even letters have assigned degree two and the orbit letters degree one. The measure-valued Magnus map remains injective after scalar extension and this invertible homogeneous change.

\paragraph{Colored pair tilings and Laurent coefficient separation.}
Suppose $g\ne1$ is in the kernel and let $d$ be its first nonzero assigned degree. A coefficient pattern is an ordered tiling of $d$ target positions by singleton odd letters and paired even letters. Each pair is colored by $(e,\alpha,k)$. For fixed target residue classes, its vanishing target coefficient gives
\[
 \sum_{\mathcal P} A_{\mathcal P}(T_i:i\text{ singleton})
 \prod_{(j,e,\alpha,k)\in\mathcal P}
 \frac{C_{e;a_j,a_{j+1}}(T_j,T_{j+1})}{D_{j,\alpha}^k}=0,
 \quad D_{j,\alpha}=(1+T_j)(1+T_{j+1})-\alpha.                \tag{G.9}
\]
All $A_{\mathcal P}$ are bounded $L$-valued measure transforms. Only the indicated first source degree contributes: a larger degree has larger target length, and every smaller degree is zero. Substitution at each finite arithmetic weight cutoff is continuous and multiplicative by integration of bounded continuous coefficient functions and uniform locally constant approximation, as in the preceding complete proofs.

There are finitely many patterns and denominators. Clearing denominators in (G.9) gives a bounded-coefficient analytic identity on the full Cartesian tails of the allowed weight grids. Successive one-variable identity theorems give an identity of power series over $L$. Its bounded-denominator ring is
\[
 R_L=\mathcal O_L[[T_1,\ldots,T_d]][1/\pi],
\]
not the ring permitting arbitrary unbounded coefficient denominators.

Choose a nonzero pattern with the maximum number of even letters. Fix its pair positions $J$ and its $(e,\alpha)$ colors. Among nonzero patterns with those positions and colors choose a pole-order vector $(k_j)$ maximal in the componentwise partial order. This is possible because there are finitely many patterns. Assign the unit target pair $(a_e,b_e)$ from (G.3) at each chosen pair and its source orbit type at each singleton.

For $j\in J$ introduce $S_j=D_{j,\alpha_j}$. The pairs are disjoint and the substitutions
\[
 T_{j+1}=\frac{\alpha_j+S_j}{1+T_j}-1
\]
give an isomorphism of formal $\mathcal O_L$ power-series rings: their constant terms are topologically nilpotent. Retain the unpaired variables and one left variable from each pair as base variables. Let $K_0$ be the fraction field of their bounded-denominator power-series ring. Equation (G.9) embeds into Laurent series over $K_0$ in the $S_j$, with finite lower bounds on all exponents.

Every other denominator is analytic in these variables. At the same pair position a different root gives $S_j+(\alpha_j-\alpha')$ with nonzero constant in $K_0$. At a different position its two entries lie in different blocks, whose remaining base variables are independent; their product is not identically its root. Its constant term is again nonzero in $K_0$. These denominators are therefore invertible in $K_0[[S_j:j\in J]]$. A different even character at a selected position contributes zero for the assigned target types, since their sum fixes $e$. Each selected $C$ remains a nonzero unit at $S=0$.

Extract the coefficient of $\prod_{j\in J}S_j^{-k_j}$. A pattern missing a selected pair or using a different root there has no negative exponent in that variable. A pattern properly containing $J$ has more even letters and zero coefficient by maximality. For patterns with the same positions and root colors, an order vector $k'$ can contribute only if $k'_j\ge k_j$ for every $j$, since all remaining series have nonnegative exponents. Componentwise maximality excludes every strict such vector. Thus only the selected pattern contributes, and its coefficient is
\[
 A_{\mathcal P}\prod_{j\in J}
 C_{e_j;a_{e_j},b_{e_j}}\big|_{S=0}\ne0.
\]
The unpaired-variable coefficient ring embeds in $K_0$, and all factors are units. This is a contradiction. When there are no even pairs, the same calculation reduces to odd-moment separation. Hence the actual rational action is injective on $H$. Cyclotomic abelianization detects any kernel outside $H$, and the rational action is induced by the pro-$p$ outer action. Theorem~\ref{thm:pole-general} follows.

\paragraph{K17-CUP-PARITY: from $p$-units to odd Galois generators.}
Let $A_F$ be the $p$-class group and
\[
 E=\mathcal O_F[1/p]^\times/(\mathcal O_F[1/p]^\times)^p.
\]
For a $p$-unit $v$, $v/c(v)$ is a global unit with all complex conjugates of absolute value one, hence a root of unity. An elementary proof is that the monic polynomials of its positive powers have bounded degree and bounded integer coefficients, so only finitely many roots occur and two powers coincide. Thus $E^-$ is exactly the one-dimensional space generated by $\mu_p$; it is nonzero because $F$ does not contain $\mu_{p^2}$.

The Kummer sequence and Vandiver give
\[
 0\longrightarrow E\longrightarrow H^1(P,\mu_p)
  \longrightarrow A_F[p]\longrightarrow0,
 \qquad E^+=H^1(P,\mu_p)^+.
\]
There is no Brauer contribution, since $F$ is totally imaginary with just one omitted finite prime, and
\[
 H^2(P,\mu_p^{\otimes2})=A_F\otimes\mu_p
\]
is a plus space. Mixed parity cups therefore vanish. The alternating cup on the one-dimensional $E^-$ vanishes because $p$ is odd. Surjectivity of the cup on all $p$-units consequently implies surjectivity from $E^+\otimes E^+$. Untwisting gives precisely the odd $H^1(P,\mathbb F_p)$ cup hypothesis of Theorem~\ref{thm:pole-general}. Under cup/relator duality, $x_a$ has character $\omega^a$, its Kummer dual has character $\omega^{1-a}$, and a relation of character $\omega^e$ corresponds to twisted target character $\omega^{2-e}$. Thus the indices satisfy $a+b\equiv e\pmod{p-1}$, as used above.

\begin{corollary}[K17-KERNEL-LT1000]\label{cor:kernel-lt1000}
For every odd prime $p<1000$, one has $\Omega_p^*=\Omega_p$.
\end{corollary}
\begin{proof}
Buhler--Harvey verify Vandiver for every prime below $163{,}577{,}356$, in particular for this entire range \cite[abstract and \S1, p.~1]{BuhlerHarvey2011}. Sharifi's Theorem~1.2 and Corollary~5.9 prove that the $p$-unit cup pairing is surjective for every odd $p<1000$ \cite[pp.~2 and~27]{Sharifi2007}. The parity lemma supplies the exact odd-generator hypothesis, so Theorem~\ref{thm:pole-general} applies.
\end{proof}

The cited computational theorems are used as published results: Sharifi's Theorem~5.7 is an explicit Magma modular-symbol computation, and the Vandiver range is the published Buhler--Harvey verification. These computations were not rerun here. The deduction imposes no separate numerical $\lambda$-invariant or simple-root condition. For an arbitrary remaining prime, the explicit sufficient conditions are Vandiver and the stated odd cup surjectivity. This argument does not settle their validity for all primes, and it does not prove integral odd generation of the image Lie algebra.

\subsection{A second explicit integral defect: weight twelve at 691}

\begin{corollary}[T-INTEGRAL691]
At $p=691$, every $K_{4k}(\mathbb Z)\{691\}$ vanishes, but the
actual integral tripod Lie algebra is not generated by its odd
Soul\'e classes. With the notation and normalization of
McCallum--Sharifi, one has
\[
 \mathfrak g_{691,12}/\mathfrak h_{691,12}
 \simeq (\mathfrak g_{691}^{\mathrm{ab}})_{12}
 \simeq\mathbb Z/691,
\]
and
\[
 [\sigma_3,\sigma_9]\equiv50[\sigma_5,\sigma_7]
       \pmod{691\mathfrak g_{691,12}}.
\]
Here $\mathfrak h$ denotes the integral Lie subalgebra generated by
the specified odd classes. The displayed coefficient $50$ belongs
to that normalization.
\end{corollary}

\begin{proof}
Vandiver at $691$ lies in the already cited verified range, so
VK-EQUIV gives the $K$-group assertion. Sharifi's cup-product
surjectivity theorem, \cite[Theorem~1.2, p.~2; Corollary~5.9,
p.~27]{Sharifi2007}, gives nontriviality in the irregular component
$(691,12)$. The same paper explicitly makes the application to
Ihara's weight-twelve relation on p.~28.
McCallum--Sharifi \cite[Theorem~9.11, p.~30]{McCallumSharifi2003}
then gives strict inclusion $\mathfrak h_{691,12}\subsetneq
\mathfrak g_{691,12}$ and the displayed congruence. Since $12$ is
the least irregular index, Theorem~9.7 identifies this quotient
with the degree-twelve abelianization. Proposition~9.10 and its
proof, pp.~29--30, identify its nonzero image with the cyclic
group of order $691$ in the stable derivation algebra. This proves
the claimed exact quotient. These are applications of the cited
results, not a new computation of the stable derivation algebra.
\end{proof}

Cup-product surjectivity alone does not assert that the defect for
every irregular prime first occurs at its least irregular index.
Theorem~9.7 additionally requires a nonzero coefficient with
$i+j=m$ and $3\le i<j$; a nonzero pairing whose least positive
representatives sum to $m+p-1$ does not meet that hypothesis.
The two explicit examples $37$ and $691$ have the required
published low-weight computations.

\subsection{What remains integral at \texorpdfstring{$2$}{2}}

Theorem~\ref{thm:pro2-kernel} and rational Deligne--Ihara do not prove surjectivity of
\[
 \operatorname{Lie}_{\mathbb Z_2}(s_3,s_5,\ldots)
      \longrightarrow\mathfrak g_2.
\]
Nonzero integral odd classes exist as explained above. Choosing odd classes whose rational images form the standard free generating system gives an injective integral free source. Equality with the actual lattice, weight by weight, is the remaining saturation question. We obtained neither a checked all-weight proof nor a counterexample for this surjectivity at $2$.

There is a useful rational normal form, but it does not remove that question. Choose a Levi subgroup containing the image of $c$. Since $chc=h^{-1}$, the element $a=\log h$ has only odd total weights. Writing $\gamma=vu$, the element $\exp(a/2)v$ is $c$-fixed. A further even unipotent conjugation can remove this factor because $u^w-1\neq0$ in $\mathbb Q_2$ for every positive weight $w$, giving a rational form $\gamma=\exp(-a/2)u$. Division by $2$ and inversion of $u^w-1$ need not preserve an integral lattice. This suggests finite-weight lattice calculations; it supplies no integral free basis by itself.

\subsection{Mixed Tate motives and the current \texorpdfstring{$\mathfrak{grt}_1$}{grt1} statement}

Deligne--Goncharov's category over $\mathbb Z$ has rational coefficients. In the attached LaTeX, \S1.6, Propositions~1.7 and~1.9, and Proposition~2.3 establish the rational Ext and free-Lie statements; \S2.4 records the noncanonical nature of free lifts \cite{DeligneGoncharov2005}. The base $\operatorname{Spec}\mathbb Z$ specifies ramification conditions. It does not replace $\mathbb Q$-coefficients by $\mathbb Z$ or $\mathbb Z_p$. Rational $\operatorname{Ext}^2$ vanishing erases torsion and is not an integral vanishing result.

The rational Grothendieck--Teichm\"uller status must also be updated from current primary sources. Hirose's October2026 manuscript, arXiv:2610.23719v1, Theorem~4 and \S1.2, gives a Hoffman spanning result over $\mathbb Z_{(2)}$ for the confluence quotient and derives rational equality of the de Rham mixed-Tate motivic Galois group with GRT \cite{Hirose2026}. We checked the statement and its stated implication, not the whole proof. This is a separate current rational result. It does not identify the arithmetic $\mathbb Z_p$-image lattice, and it does not invalidate the counterexample at $37$. An integral associator lattice, a $\mathbb Z_{(2)}$-span of motivic periods, and the lattice of an arithmetic pro-$p$ Galois image are different objects unless a comparison theorem is supplied.

The supplied round-two war-game file mentions B3 and S9 but does not reproduce their original round-one rows. We have therefore used their subject as specified in the prompt and have not assigned an invented exact locator to missing text.

% Bibliography keys needed in this fragment and pro2_section.tex:
% Sharifi2002: Relationships between conjectures on the structure of pro-p Galois groups unramified outside p; https://www.math.ucla.edu/~sharifi/publ/galstr.pdf
% Brown2012: Mixed Tate motives over Z, Annals175(2012)949-976; https://annals.math.princeton.edu/wp-content/uploads/annals-v175-n2-p10-p.pdf
% Brown2019: From the Deligne-Ihara conjecture to multiple modular values; https://arxiv.org/pdf/1904.00179
% HainMatsumoto2001: Weighted completion of Galois groups and Galois actions on the fundamental group of P1-{0,1,infinity}; https://arxiv.org/pdf/math/0006158
% GrasJaulent2008: Note on 2-rational fields; https://arxiv.org/pdf/0805.1168
% MinacPasiniQuadrelliTan2020: Koszul algebras and quadratic duals in Galois cohomology; https://arxiv.org/pdf/1808.01695
% IshiiKernel2026: On the kernels of the pro-p outer Galois representations associated to once-punctured CM elliptic curves; arXiv2312.04196v3
% IshiiKummer2024: On Kummer characters arising from Galois actions on the pro-p fundamental groups of once-punctured CM elliptic curves; arXiv2312.04175v2
% McCallumSharifi2003: A cup product in the Galois cohomology of number fields; https://www.math.ucla.edu/~sharifi/publ/pairing.pdf
% BuhlerHarvey2011: Irregular primes to 163 million; arXiv0912.2121
% Iwasawa1981: Riemann-Hurwitz formula and p-adic Galois representations for number fields, Tohoku33(1981)263-288; https://www.jstage.jst.go.jp/article/tmj1949/33/2/33_2_263/_pdf/-char/en
% SharifiIwasawaNotes: Iwasawa Theory, Chapter3; https://www.math.ucla.edu/~sharifi/notes/iwasawa-ch03.html
% DeligneGoncharov2005: Groupes fondamentaux motiviques de Tate mixte; attached LaTeX Sec1ccc/Sec2ccc.
% Hirose2026: Confluence relations, Hoffman basis over Z_(2), and the Deligne-Drinfeld conjecture; arXiv2610.23719v1.

\subsection{Exact integral comparison and a finite calculation at 2}

\begin{proposition}[K17-OUT-SAT]\label{prop:outer-sat}
For a finitely generated free pro-$p$ group $\Pi$ of rank at least
two, every positive graded quotient of its lower-central
outer-automorphism filtration is a torsion-free $\mathbb Z_p$-module.
For an arithmetic outer action covered by Hain--Matsumoto
Proposition~8.2, the integral pullback filtration and its rational
unipotent pullback therefore agree exactly in positive degrees.
This includes $p=2$.
\end{proposition}

\begin{proof}
Write $\Gamma_n$ for the closed lower central series of $\Pi$,
$L_n=\Gamma_n/\Gamma_{n+1}$ and $V=L_1$. The graded Lie algebra
$L$ is the free $\mathbb Z_p$-Lie algebra on $V$, for every prime
$p$, by Koch--Kukkuk--Labute \cite[Theorem~2.1, pp.~48--49;
author PDF pp.~4--5]{KochKukkukLabute1998}.
Put
\[
 A_n=\ker\bigl(\operatorname{Aut}\Pi\longrightarrow
                    \operatorname{Aut}(\Pi/\Gamma_{n+1})\bigr).
\]
The map sending $\alpha$ to
$x\mapsto\alpha(x)x^{-1}\bmod\Gamma_{n+2}$ is
$\mathbb Z_p$-linear in $x\in V$, additive in $\alpha$, and has
kernel $A_{n+1}$. These assertions follow from the commutator
identities and the fact that $\alpha$ acts trivially modulo
$\Gamma_{n+1}$. Thus it embeds $A_n/A_{n+1}$ in
$\operatorname{Hom}_{\mathbb Z_p}(V,L_{n+1})$.

Every outer automorphism trivial modulo $\Gamma_{n+1}$ has an
$A_n$ representative: lift its quotient inner conjugator to
$\Pi$ and remove it. An inner conjugation by $g$ belongs to
$A_n$ precisely when $g\in\Gamma_n$. Indeed a first nonzero
degree of $g$ below $n$ would centralize all free Lie generators.
The center of a free Lie algebra of rank at least two over any
field is zero in positive degree. To see this, embed it in the
free associative algebra. A polynomial commuting with a free
letter $X$ contains only powers of $X$, by comparing word
coefficients and moving initial $X$'s to the right. Commutation
with another independent letter then makes it scalar, hence zero
in positive degree. The characteristic-zero conclusion also
holds over $\mathbb Z_p$ by tensoring.

It follows that the positive outer graded quotient embeds in
\[
 \operatorname{Hom}_{\mathbb Z_p}(V,L_{n+1})/\operatorname{ad}(L_n).
\]
The ambiguity is exactly the inner degree-$n$ map; if a Johnson
class is $\operatorname{ad}(z)$, an integral lift of $z$ removes
it and puts the representative in $A_{n+1}$. The displayed
quotient is torsion-free. In fact $\operatorname{ad}$ is injective
modulo $p$, since the free Lie algebra over $\mathbb F_p$ has zero
center. Thus $\operatorname{ad}(z)$ divisible by $p$ forces $z$
divisible by $p$, so its image is primitive in the finite free
target. This also proves torsion-freeness for every arithmetic
subgroup with the induced filtration.

Let $I^n\subseteq J^n$ be the integral and rational pullbacks.
They agree at $n=1$, since identity on the free integral
abelianization is equivalent to identity after tensoring with
$\mathbb Q_p$. If $I^n=J^n$, Hain--Matsumoto
\cite[Proposition~8.2, p.~25]{HainMatsumoto2001} gives a finite
kernel on $\operatorname{gr}_I^n\to\operatorname{gr}_J^n$.
The source is torsion-free, so the kernel is zero. Every element
of $J^{n+1}\subset I^n$ maps to that kernel and is in $I^{n+1}$.
Induction proves the asserted exact equality.
\end{proof}

\begin{theorem}[K17-2-FINITE18]\label{thm:2-finite18}
There are primitive integral odd classes $\sigma_m$ in the actual
tripod image at $2$ such that, in every arithmetic weight at most
$18$, their Hall brackets generate the entire integral graded
Lie algebra. This is a finite-range result; the all-weight
integral generation assertion remains open in this run.
\end{theorem}

\begin{proof}
We give the mathematical reduction and the exact finite
certificate. Use the checked presentation $C_2*\mathbb Z_2$,
choose $\chi(\gamma)=5$, and put $h=c\gamma c\gamma^{-1}$.
After placing $c$ in a rational Levi torus, $a=\log h$ has only
odd weights because $chc=h^{-1}$. Write $\gamma=v u$ with
$u=5$. The element $z=\exp(a/2)v$ is $c$-fixed. A $c$-fixed
conjugation removes $z$ successively in positive even weights:
the equation in weight $w$ divides by $5^w-1\ne0$ in
$\mathbb Q_2$. We obtain
\[
 h=\exp(a),\qquad \gamma=\exp(-a/2)u.
\]
Every odd homogeneous component of $a$ has nonzero
indecomposable coordinate, by density and the compact orbit
generation of the cyclotomic kernel. Brown's rational freeness
therefore permits choosing these components as the free
generators $s_m$, so that
\[
 a=\sum_{m\ge3,\ m\ {
m odd}}s_m,\qquad
 D=\operatorname{Ad}(\exp(-a/2))\circ\tau,
 \quad\tau(s_m)=5^m s_m.
\]
All changes here transport the actual compact image and its
lattice; they do not declare the $s_m$ to be integral generators.
By Proposition~\ref{prop:outer-sat}, its transported weight
filtration is exactly the original integral filtration.

At cutoff $W$, work in the finite-dimensional free associative
algebra over $\mathbb Q$ truncated by arithmetic weight. The
finite exponential and logarithm give exact products and
commutators, retaining every denominator. The desired compact
group $H_W$ is the closed group generated by the continuous
$D$-orbit of $\exp(a)$.

Order a Hall--Lyndon basis by increasing weight, putting the
primitive letter first within each odd weight. Store group
generators by their logarithms. If a candidate has first
coordinate $b$ and the current pivot has first coordinate $c$,
eliminate it by the corresponding group power whenever
$v_2(b/c)\ge0$. If its valuation is smaller, replace the pivot
and reduce the old pivot in turn. Normalize only by unit
$\mathbb Z_2$-powers. Add conjugates under $D$ and pair
commutators until the output stabilizes. Every constructed
element belongs to $H_W$.

Here is why the finite checks certify the whole compact group.
For a triangular final list $g_1,\ldots,g_N$, put
\[
 C_i=g_i^{\mathbb Z_2}g_{i+1}^{\mathbb Z_2}\cdots
                         g_N^{\mathbb Z_2}.
\]
Every commutator has greater weight than both inputs. Its checked
zero reduction expresses it in the later ordered factors.
Descending induction makes each $C_i$ a subgroup: $g_i$
normalizes the already closed subgroup $C_{i+1}$, first for
nonnegative integral powers and then for all $\mathbb Z_2$
powers by continuity. In particular
$g_i^{2^r-1}\to g_i^{-1}$ supplies inverse conjugation.
Each $C_i$ is compact by its continuous parametrization.
Triangular first coordinates make that parametrization
injective and successive reduction an exact membership test.

The final certificate verifies that $h$, every $D(g_i)$,
every $D^{-1}(g_i)$ and every potentially nonzero pair
commutator reduce to the identity. Consequently $C_1$ contains
the whole compact orbit and $H_W\subseteq C_1$. The construction
gives the reverse inclusion. Hence $C_1=H_W$.

Vanishing of coordinates below weight $n$ is equivalent to
vanishing of all preceding product exponents. Its complete
weight-$n$ lattice is thus the $\mathbb Z_2$-span of the leading
weight-$n$ pivot vectors. Select the primitive odd pivots and
compute all their Hall brackets. They are contained in this
lattice. The quotient of the two determinants measures its
index; equality of their $2$-adic valuations proves equality
of the lattices.

The exact rational computation at $W=18$ gives the following
table. Columns $d_n$ and $d_n^{\rm odd}$ are determinant
valuations in the stated rational Hall coordinates.
\[
\begin{array}{c|r|r|r@{\qquad}c|r|r|r}
n&\mathrm{rank}&d_n&d_n^{\rm odd}&
n&\mathrm{rank}&d_n&d_n^{\rm odd}\\\hline
3&1&0&0&5&1&3&3\\
7&1&7&7&8&1&3&3\\
9&1&10&10&10&1&7&7\\
11&2&18&18&12&2&20&20\\
13&3&31&31&14&3&31&31\\
15&4&52&52&16&5&66&66\\
17&7&96&96&18&8&114&114
\end{array}
\]
All other positive weights at most $18$ have zero rational
space and hence zero integral graded group. Every listed index
is one. The complete rational vectors, arithmetic operations,
closure checks and output are in
\path{evidence/work/pro2_integral_malcev.py} and
\path{evidence/work/pro2_integral_malcev_extended.json}.
An independent rerun through $14$ and independent two-way
integral change-of-basis checks through $18$ are recorded in
the accompanying audit JSON files. The algorithm and the
collection and comparison proofs were read in full. This is
an exact reproducible computation, not a formal-verifier
certificate.
\end{proof}

The observed primitive odd valuation at weight $m=2r+3$ is
$3r+v_2(r!)$. Its abelian lower bound is exact:
\[
 v_2\!\left(\prod_{j=0}^{r-1}
 (5^{2r+3}-5^{2j+3})\right)=3r+v_2(r!),
\]
since $v_2(5^{2s}-1)=3+v_2(s)$. Successive division by the
monic linear factors describes the preceding-moment kernel
in the Iwasawa algebra. The unproved step in an all-weight
argument is integral correction of every lower decomposable
weight and stability of the resulting Hall family under $D$.
The scalar Newton determinant does not alone prove that step.

\section{Cote 149, scans 79--85: Kummer theory and geometric tangential splittings}
\label{sec:archive-kummer}

\paragraph{Evidence and notation.}
All seven supplied TeX transcriptions were read in full.  The original
manuscript images were not supplied, so the editorial conclusions below are
mathematical readings, not verified readings of handwritten glyphs.  Write
\(D\) for the boundary of a smooth proper relative curve \(X/S\), put
\(U=X\setminus D\), and let \(g_i:S\longrightarrow D\) be a boundary section.
This avoids the manuscript's reuse of \(T\), which we reserve for a Tate module.
Write \(E_S=\pi_1(S)\), with geometric base points suppressed.

The passage studies a proposed cusp-normalizer extension
\begin{equation}\label{eq:arc-normalizer}
 1\longrightarrow T\longrightarrow N(L_i)
   \longrightarrow E_S\longrightarrow 1,
 \qquad T=\widehat{\mathbb Z}(1),
\end{equation}
and its obstruction class \(c_i\).  The geometric interpretation proposed in
scan 81 is the Kummer first Chern class of the conormal or normal line of
\(g_i(S)\).  The author explicitly leaves that sign convention unsettled.
Scan 85 then asks for an intrinsic profinite characterization of the
\emph{geometric} splittings inside the set of all splittings.  The obstruction
to the existence of a splitting and the recognition of that subset are
different questions.

In the calculations below, \(\mathbb Z_\ell(1)\) means the inverse system
\((\mu_{\ell^r})_r\), and \(H^j_{\mathrm{cont}}\) means continuous
cohomology of that inverse system.  We do not silently replace it by
ordinary cohomology of the inverse-limit sheaf on the small etale site.
For the full Tate module, the corresponding index set is all positive
integers ordered by divisibility.

\paragraph{F149-KUM: exact limiting Kummer sequences (proved).}
Let \(S\) be connected, normal and noetherian, and suppose that \(\ell\)
is invertible on \(S\).  Set
\[
 A=\Gamma(S,\mathcal O_S)^\times,\qquad
 P=\operatorname{Pic}(S),\qquad
 B=H^2_{\mathrm{et}}(S,\mathbb G_m).
\]
For an abelian group \(M\), put
\[
 M^\wedge_\ell=\varprojlim_r M/\ell^rM,
 \qquad
 T_\ell M=\varprojlim_r M[\ell^r],
\]
where the second transition map is multiplication by \(\ell\);
for \(A\), written multiplicatively, the first quotient is
\(A/A^{\ell^r}\).  There are natural short exact sequences
\begin{align}
 0&\longrightarrow A^\wedge_\ell
   \longrightarrow \varprojlim_r H^1_{\mathrm{et}}(S,\mu_{\ell^r})
   \longrightarrow T_\ell P\longrightarrow 0,\label{eq:arc-kum-h1}\\
 0&\longrightarrow P^\wedge_\ell
   \longrightarrow \varprojlim_r H^2_{\mathrm{et}}(S,\mu_{\ell^r})
   \longrightarrow T_\ell B\longrightarrow 0.\label{eq:arc-kum-h2}
\end{align}
Moreover,
\[
 H^1_{\mathrm{cont}}(S,\mathbb Z_\ell(1))
 \simeq \varprojlim_r H^1_{\mathrm{et}}(S,\mu_{\ell^r}).
\]
If \(P[\ell^\infty]\) has bounded exponent, then \(T_\ell P=0\), the
system \(H^1_{\mathrm{et}}(S,\mu_{\ell^r})\) is Mittag--Leffler, and
\begin{align*}
 H^1_{\mathrm{cont}}(S,\mathbb Z_\ell(1))&\simeq A^\wedge_\ell,\\
 H^2_{\mathrm{cont}}(S,\mathbb Z_\ell(1))
 &\simeq \varprojlim_r H^2_{\mathrm{et}}(S,\mu_{\ell^r}).
\end{align*}
In particular the extra hypothesis holds when \(P\) is finitely generated.

\paragraph{Proof.}
The etale Kummer sequence and Hilbert 90 give
\begin{align*}
 0&\longrightarrow A/A^{\ell^r}
   \longrightarrow H^1_{\mathrm{et}}(S,\mu_{\ell^r})
   \longrightarrow P[\ell^r]\longrightarrow 0,\\
 0&\longrightarrow P/\ell^rP
   \longrightarrow H^2_{\mathrm{et}}(S,\mu_{\ell^r})
   \longrightarrow B[\ell^r]\longrightarrow 0.
\end{align*}
The transition from level \(r+1\) to level \(r\) is induced by raising
roots of unity to the \(\ell\)-th power.  In the Kummer diagram the first
copy of \(\mathbb G_m\) maps by that power and the second by the identity.
Consequently the transitions on the left are quotient maps, and the
transitions on the right are multiplication by \(\ell\).  Both left
systems are surjective, so their first derived inverse limits vanish.
Taking inverse limits proves
\eqref{eq:arc-kum-h1}--\eqref{eq:arc-kum-h2}, including the terminal zeros.

If \(\ell^bP[\ell^\infty]=0\), the transition
\(P[\ell^m]\longrightarrow P[\ell^r]\) is zero whenever \(m\geq r+b\).
The image of \(H^1(S,\mu_{\ell^m})\) at level \(r\) is then contained in
the unit term \(A/A^{\ell^r}\).  It contains that entire term because
the unit quotient transitions are surjective.  These stabilized images
prove the Mittag--Leffler assertion and also give \(T_\ell P=0\).
For a countable inverse system, the continuous-cohomology comparison is
\begin{equation}\label{eq:arc-jannsen}
 0\longrightarrow \varprojlim\nolimits^1_r
 H^{j-1}_{\mathrm{et}}(S,\mu_{\ell^r})
 \longrightarrow H^j_{\mathrm{cont}}(S,\mathbb Z_\ell(1))
 \longrightarrow \varprojlim_r
 H^j_{\mathrm{et}}(S,\mu_{\ell^r})\longrightarrow 0.
\end{equation}
The groups \(H^0(S,\mu_{\ell^r})\) are finite: a section is determined
at the generic point of the connected normal scheme and satisfies a
polynomial of degree \(\ell^r\).  Their inverse system is therefore
Mittag--Leffler.  This removes the degree-one derived-limit term without
a Picard hypothesis.  The preceding calculation removes the
degree-two term under the bounded-torsion hypothesis.  This proves all
the assertions.

The finite-level inputs are \emph{SGA 4}~\cite{SGA4}, Expose IX,
Theorie de Kummer 3.2 and Theoreme 3.3; the primary French reedition was
checked at printed pp.~14--15.  Formula \eqref{eq:arc-jannsen} is Jannsen~\cite{Jannsen1988},
\emph{Continuous etale cohomology}, \emph{Math.\ Ann.} 280 (1988),
equation (3.1), p.~216; the inverse-system Kummer diagram is
equation (3.27), p.~222.  Both locators were checked in the primary text.

\paragraph{Full completion and its precise scope.}
If \(S\) is a scheme over \(\mathbb Q\), all positive integers are
invertible.  The same proof, on the cofinal sequence \(n=r!\), gives
\begin{align*}
 0&\longrightarrow A^\wedge
 \longrightarrow H^1_{\mathrm{cont}}(S,\widehat{\mathbb Z}(1))
 \longrightarrow \varprojlim_n P[n]\longrightarrow 0,\\
 0&\longrightarrow P^\wedge
 \longrightarrow \varprojlim_n H^2_{\mathrm{et}}(S,\mu_n)
 \longrightarrow \varprojlim_n B[n]\longrightarrow 0,
\end{align*}
where \(A^\wedge=\varprojlim_n A/A^n\) and
\(P^\wedge=\varprojlim_n P/nP\).  If the torsion of \(P\) has bounded
exponent, the same stabilization argument applies to the factorial
system and identifies the middle degree-two group with continuous
cohomology.  This is the pro-\(\mathbb N\) completion.  It need not be
a compact profinite group when the quotients \(A/A^n\) are infinite.

Thus the uncertain terminal arrow in scan 82, formula (21), can be
resolved mathematically as an arrow to zero with the usual transitions.
The \(K(\pi,1)\) hypothesis in scan 80 does not by itself justify dropping
every derived inverse limit.  In the finite-generation setting the
argument above supplies the missing justification.  These observations
are not a claim to have visually read the uncertain arrow.

The requirement that the relevant integers be invertible is substantive.
For example, over \(k=\mathbb F_p(t)\), the element \(t\) acquires no
\(p\)-th root in any separable algebraic extension: such a root would
generate the nontrivial purely inseparable extension defined by
\(X^p-t\).  Hence \(x\mapsto x^p\) on \(\mathbb G_m\) is not an epimorphism
on the etale site of \(k\).  Kummer theory for noninvertible exponents
uses the flat topology; its finite flat torsors cannot simply be
reinterpreted as etale fundamental-group torsors.  In particular, the
example over \(\operatorname{Spec}\mathbb Z[1/\ell]\) below is a
\emph{pro-\(\ell\)} example, not an unqualified full-\(\widehat{\mathbb Z}\)
etale calculation on that mixed-characteristic scheme.

\paragraph{F149-SPLIT: the two torsors (proved).}
Fix a split extension \(1\to T\to E\to G\to1\), with \(T\) abelian,
and choose a continuous splitting \(s_0\).  Every other splitting is
uniquely \(s(g)=a(g)s_0(g)\), where
\[
 a(gh)=a(g)\,{}^ga(h).
\]
This is exactly the continuous cocycle equation.  Conjugating a
splitting by an element of \(T\) changes \(a\) by a coboundary.
Consequently the \emph{conjugacy classes} of splittings form a torsor
under \(H^1_{\mathrm{cont}}(G,T)\).  Raw splittings are described by
\(Z^1_{\mathrm{cont}}\), not by \(H^1_{\mathrm{cont}}\).
For the prime-invertible systems used here, the degree-one group and
etale descriptions agree: finite \(\mu_n\)-torsors are finite etale,
and their degree-zero inverse system is Mittag--Leffler as above.
No degree-two \(K(\pi,1)\) comparison is needed for this assertion.

If a line bundle \(\mathcal L\) on \(S\) is trivial, its global bases form
a torsor under \(A=\Gamma(S,\mathcal O_S)^\times\).  Indeed, two bases
\(b,b'\) have a unique ratio \(u\in A\), and every \(u\in A\) occurs.
In a local tame parameter model, replacing \(t\) by \(ut\) changes the
root tower by roots of \(u\).  Its descent cocycle is precisely the
Kummer cocycle of \(u\).  Thus the map from bases to local splitting
classes is equivariant for
\[
 A\longrightarrow H^1_{\mathrm{cont}}(S,T).
\]
The same assertion for the proposed global normalizer follows when the
relative comparison specified below has been constructed.

Higher-order parameter changes introduce no additional finite-level
class: if \(n\) is invertible in a ring \(R\), the map
\(v\mapsto v^n\) is a bijection of \(1+tR[[t]]\).
To prove this, solve successively for the coefficients of
\(v=1+\sum_{j\geq1}v_jt^j\).  At stage \(j\) the new coefficient appears
as \(nv_j\), while every other term is determined by earlier
coefficients.  The invertibility of \(n\) gives existence and uniqueness
at every stage.  A parameter change \(t\mapsto ut+O(t^2)\) therefore has
the same Kummer ambiguity as its leading unit.

\paragraph{F149-FULLINJ: injectivity of the full unit completion (proved).}
If \(S\) is an integral scheme whose function field \(F\) is finitely
generated over \(\mathbb Q\), then
\[
 \Gamma(S,\mathcal O_S)^\times
 \longrightarrow
 \varprojlim_n\Gamma(S,\mathcal O_S)^\times/
       \bigl(\Gamma(S,\mathcal O_S)^\times\bigr)^n
\]
is injective.

\paragraph{Proof.}
It suffices to prove \(\bigcap_n(F^\times)^n=\{1\}\).  Let \(k\) be the
algebraic closure of \(\mathbb Q\) in \(F\); it is a number field.
Choose a normal projective \(k\)-model of \(F\).  If \(u=v_n^n\) for
every \(n\), the order of \(u\) at every prime divisor is divisible by
every positive integer, hence is zero.  A normal noetherian domain is
the intersection, in its fraction field, of its local rings at
height-one primes.  Applying this to affine opens, both \(u\) and
\(u^{-1}\) are global regular functions on the model.  Properness then
puts \(u\) in \(k^\times\); the chosen constant field is algebraically
closed in \(F\).  Each \(v_n\), being algebraic over \(k\), also lies in
\(k^\times\).  Its finite-prime valuations show that \(u\) and every
\(v_n\) are units of the ring of integers of \(k\).  Dirichlet's unit
theorem makes this unit group finitely generated.  The intersection
of all \(n\)-multiple subgroups of a finitely generated abelian group
is zero, by its decomposition into a free group and a finite group.
Thus \(u=1\), as required.

This proves the full-completion assertion used in scan 85 under the
stated arithmetic function-field hypothesis.  It does not prove the
analogous assertion for every single prime.  For every odd \(\ell\),
\(-1=(-1)^{\ell^r}\) for all \(r\), so \(-1\) maps to zero in the
\(\ell\)-adic completion.  The full and one-prime statements must not
be interchanged.

\paragraph{F149-UNIT: a nongeometric local splitting class for every prime
(proved).}
Let \(S_\ell=\operatorname{Spec}\mathbb Z[1/\ell]\).  Then
\[
 H^1_{\mathrm{cont}}(S_\ell,\mathbb Z_\ell(1))
 \simeq
 \begin{cases}
 \mathbb Z_\ell,&\ell\ne2,\\
 \mathbb Z/2\mathbb Z\oplus\mathbb Z_2,&\ell=2.
 \end{cases}
\]
In the local model of a standard tripod cusp with a fixed parameter,
there is a local splitting class whose valuation coordinate is
\[
 \alpha_\ell=\frac{1}{1-\ell^2}\in\mathbb Z_\ell\setminus\mathbb Z.
 \tag{F149-UNIT}
\]
It is not induced by any global unit or global tangent basis, despite
\(\operatorname{Pic}(S_\ell)=0\).

\paragraph{Proof.}
The localization \(\mathbb Z[1/\ell]\) is a principal ideal domain and
has unit group \(\{\pm1\}\times\ell^{\mathbb Z}\).  Thus its Picard group
is zero, and F149-KUM identifies the displayed cohomology with the
completion of this unit group.  For odd \(\ell\), raising to
\(\ell^r\) is an automorphism of \(\{\pm1\}\); this factor disappears
in every quotient.  For \(\ell=2\) it remains a copy of
\(\mathbb Z/2\mathbb Z\), with identity transitions.  In either case
the valuation quotients are \(\mathbb Z/\ell^r\mathbb Z\), with their
usual transitions.  This proves the group formula.

The denominator \(1-\ell^2\) is a unit in \(\mathbb Z_\ell\), while
its absolute value as a reduced integer denominator is
\(\ell^2-1>1\).  Hence \(\alpha_\ell\) is \(\ell\)-adically integral
but is not an integer.  Choose integers \(a_r\) with
\(a_r\equiv\alpha_\ell\pmod{\ell^r}\).  The Kummer classes of
\(\ell^{a_r}\) are compatible, since
\(a_{r+1}-a_r\) is divisible by \(\ell^r\), and they do not depend
on the choices.  If their inverse-limit class came from a global
unit \(\pm\ell^a\), its valuation coordinate would give
\(a=\alpha_\ell\) in \(\mathbb Z_\ell\), a contradiction.

The chosen standard parameter supplies the split local tame extension.
Its splitting classes are the torsor in F149-SPLIT; changing a global
tangent basis is exactly the Kummer action of a global unit.  The
constructed class lies outside that image.  The conclusion concerns
the local splitting torsor.  It neither assumes a general relative
normalizer theorem nor proves that a richer anabelian structure cannot
recognize the geometric subset.

\paragraph{F149-TORSOR: the completion mark in scan 85.}
The group acting on \emph{bases} is the uncompleted group
\(\mathbb G_m(S)\).  Scan 82, line 4, states this correctly.  When the
Kummer map is injective, the geometric image is a sub-torsor under
that uncompleted group, whereas all splitting classes are acted on by
\(H^1(S,T)\).  F149-UNIT exhibits explicitly why replacing the former
acting group by its completion changes the statement.  Therefore the
flagged uncertain \(\wedge\) in scan 85, line 8, should mathematically
be absent if the sentence refers to the geometric bases and their
image.  Here and below an uncertainty command in the supplied source
is a source annotation, not an assertion that the original glyph was
inspected.

\paragraph{F149-CHERN-LOCAL: the finite-level line-bundle obstruction
(proved locally).}
Let \(\mathcal L=J_i/J_i^2\) be the conormal line, and let \(n\) be
invertible on \(S\).  The local root-coordinate gluing obstruction is
the Kummer boundary
\[
 \delta_n(\mathcal L)\in H^2_{\mathrm{et}}(S,\mu_n),
\]
up to inversion if the normal rather than conormal convention is used.

\paragraph{Proof.}
Choose local frames and formal parameters \(t_a\), with leading
transition units \(u_{ab}\).  The \(u_{ab}\) represent the line bundle.
On an etale refinement of the overlaps, choose
\(v_{ab}\) with \(v_{ab}^n=u_{ab}\); such refinements exist because
\(n\) is invertible.  The failure of the root-coordinate identifications
to compose on triple overlaps is
\[
 v_{ab}v_{bc}v_{ca}\in\mu_n.
\]
This is the defining boundary cocycle for the Kummer connecting map.
If necessary one uses the corresponding refinement of the Cech
hypercover; no assertion that one fixed Zariski cover computes all of
\(H^2\) is required.  The higher-order formal unit factors have unique
\(n\)-th roots by the coefficient recursion in F149-SPLIT, and do not
change this obstruction.  Dualizing \(\mathcal L\) inverts the
transition units, hence negates its cohomology class.

The local tame-cover input is \emph{SGA 1}~\cite{SGA1}, Expose XIII,
Proposition 5.2 and Corollaire 5.3, reedition pp.~315--316
(original marginal pp.~428--430).  The relative statements are
Proposition 5.5 and Corollaire 5.6.  The discrete-valuation
Abhyankar lemma is Expose X, Lemme 3.6.  The local statements and
their Kummer description were checked in the primary reedition.

\paragraph{The remaining relative comparison and recognition gates.}
The calculation just made does not identify the complete formal-tube
extension with \(N(L_i)\) in every generality suggested by scan 79.
That identification must specify the homotopy-exactness hypotheses,
the injection and normalizer of geometric inertia, and the descent
from the local tube.  A group-cohomology comparison with
\(H^2_{\mathrm{cont}}(S,T)\) must also be supplied where it is used.
The excerpt itself marks part of that setup as requiring verification.
We leave this general relative comparison as an explicit gate.

Once that comparison identifies the obstruction of
\eqref{eq:arc-normalizer} with the full Kummer Chern class, and once
\(\operatorname{Pic}(S)\) is finitely generated, the line bundle is
trivial if and only if that obstruction vanishes.  Indeed,
\(\bigcap_n n\operatorname{Pic}(S)=0\) by the structure theorem for
finitely generated abelian groups, and F149-KUM injects its completion
into the relevant degree-two group.  For a single prime, prime-to-\(\ell\)
Picard torsion is lost, so this full-completion conclusion cannot
simply be repeated.

A basis of the conormal line is the precise first-order datum needed
for a tangent vector.  A globally chosen formal parameter or a single
global regular function defining the divisor is a stronger datum in
general; the word \emph{uniformizer} in the excerpt should not be used
to assert such a global lifting without its hypotheses.

Over a field of characteristic zero the local splitting torsor is
already established in Stix~\cite{StixCuspidal2012}, \emph{On cuspidal sections of algebraic
fundamental groups}, author version dated 24 February 2012,
Proposition 11, p.~7.  Proposition 16, p.~8, identifies inertia
normalizers for a hyperbolic curve.  Theorem 17, p.~8, and Corollary 19,
p.~9, address injection and abundance of cusp packets under their
stated section-injectivity hypotheses.  These checked results concern
the local packet and all its section classes; they do not select the
uncompleted-unit image inside it.  Corollary 32, p.~12, recognizes
cuspidality under the non-Tate-character hypothesis, which is again
different from the requested geometric tangent subset.

\paragraph{Mathematical corrections versus unresolved readings.}
There are four evidence-based editorial entries:
\begin{enumerate}
 \item Scan 85, line 8: the acting group for global tangent bases is
 \(\mathbb G_m(S)\), as proved above; the uncertain completion mark
 should be absent in that interpretation.
 \item Scan 83, line 5: the generic points of irreducible divisors of
 \(S\) are codimension-one points of \(S\), not of the relative curve
 \(X\).  The subsequent normalization of \(S\) in its function field
 confirms this referent.  This is a mathematical correction; without
 the image it may be an author slip rather than a transcription error.
 \item Scan 82, formula (21): its terminal map is surjective and may
 end in zero by F149-KUM.  The handwritten choice has not been checked.
 \item Scan 80, line 2: a literal inference from surjectivity of a
 normalizer to the existence of a splitting is invalid.  The exact
 sequence
 \[
 0\longrightarrow\mathbb Z_\ell
 \xrightarrow{\;\ell\;}\mathbb Z_\ell
 \longrightarrow\mathbb Z/\ell\mathbb Z\longrightarrow0
 \]
 is nonsplit: a section would send a nonzero element of order \(\ell\)
 to an element of order \(\ell\) in the torsion-free group
 \(\mathbb Z_\ell\).  The following pages themselves study a splitting
 obstruction.  This proves that the unconditional implication cannot
 be right; it does not determine the replacement French wording.
\end{enumerate}

No correction of the normal/conormal sign is justified, because the
author expressly flags that convention.  Nor is the marginal extension
direction in scan 79 a mathematically compelled transcription error.
Once two genuine exact extensions have the same kernel and quotient,
a morphism inducing the identity on both is automatically an
isomorphism: its kernel is contained in the common kernel, where it
is injective, and a lift of the quotient followed by correction in the
kernel proves surjectivity.  The marginal passage discusses a setup
before all those exactness hypotheses have been established.

The Hartogs step in scan 83 needs normality of \(S\), rather than
merely normality of its open subset \(S'\).  Since the surrounding
setup already makes \(S\) smooth, the occurrence of \(S'\) does not
refute the argument and is not promoted to a required emendation.
The standard extension criterion is \emph{EGA}~\cite{EGAIV2}
\(\mathrm{IV}_2\), Theoreme 5.10.5, pp.~115--116, combined with
the normality criterion 5.8.6, p.~108, and Proposition 5.9.8 and
Definition 5.9.9, pp.~111--112.  These precise primary locators have
now been checked.  Normality gives Serre's \(S_2\) condition, so
the depth of \(\mathcal O_S\) is at least two at every point of a
closed complement of codimension at least two.  Theorem 5.10.5
then says that \(\mathcal O_S\) is closed with respect to that
complement; Proposition 5.9.8 identifies this with
\(\mathcal O_S\simeq j_*\mathcal O_{S'}\).
The equality of sections used in scan 83 follows.

\paragraph{Relation to K17.}
The archive uses Tate twist one.  Its degree-one ambiguity is a
completion of units, and its degree-two obstruction is a first Chern
class of a line bundle.  In algebraic \(K\)-theory, these are connected
with units in \(K_1\) and with the determinant or first
gamma-filtration part of \(K_0\); \(\operatorname{Pic}(S)\) should
not be inserted into \(K_0(S)\) as an additive subgroup without
qualification.

The negative-twist module
\((\mathbb Q_\ell/\mathbb Z_\ell)(-e)\), its strict local conditions,
and a duality with \(H^2(S,\mathbb Z_\ell(e+1))\) belong to a
different calculation.  A comparison must specify that change of
twist, the cohomological degree, the local conditions and the duality
map.  The Kummer sequences alone supply none of those identifications.
In particular, a bound or even a vanishing assertion for higher
\(K_{4k}(\mathbb Z)\) does not by itself select global units inside
their completion.  F149-UNIT leaves that distinction explicit even
when the Picard group is zero.

The supplied SA06 E.4 intake consists of status summaries, several
expressly marked as not read in full by its author.  It has not been
used as a proof of an all-twist cohomology table.  The relevant
war-game entry for scan 85 is S24: intrinsic recognition of the
geometric tangential subset.  The local calculations here settle the
cohomological and editorial distinctions, while the stated general
relative normalizer comparison and that intrinsic recognition problem
remain separate gates.

\section{Schemes, weights, and the limits of the comparison}
\label{sec:framework}

All arithmetic cohomology over the rational integers in this paper lives on
$U_\ell=\operatorname{Spec}\mathbb Z[1/\ell]$. The Tate lattice
$\mathbb Z_\ell(r)$ is the inverse system of finite etale sheaves
$\mu_{\ell^n}^{\otimes r}$, interpreted by continuous cohomology.
The discrete sheaf $\mathbb Q_\ell/\mathbb Z_\ell(-e)$ is its corresponding
divisible torsion coefficient object. These coefficient conventions,
the local conditions, and the compact/discrete duality are part of the
statements; replacing them by rational cohomology loses the finite groups
being bounded.

The weight convention is Deligne's \emph{Weil II}, Definitions 1.2.1--1.2.2,
p.~153, and 1.2.5(iv), p.~154 \cite{DeligneWeilII}: $\mathbb Q_\ell(1)$ has weight $-2$.
Thus the rational Tate object $\mathbb Q_\ell(-e)$ has weight $2e$,
and $\mathbb Q_\ell(e+1)$ has weight $-2(e+1)$. These statements refer
to rational sheaves and geometric Frobenius. They do not assign a
Frobenius purity theorem to a torsion module or prove its vanishing.
The finite $K$-groups at issue have zero rationalization even when their
integral vanishing is unknown. This elementary observation explains why
a weight or sign argument requires the additional integral congruence
and duality statements reconstructed above.

For the archive passage, \emph{SGA 4}, IX.3.2--3.3 supplies the
prime-invertible Kummer sequence and Hilbert 90. The first Chern class
is the Kummer boundary of a line bundle; \emph{SGA 5}, VII,
Proposition 3.4(i), identifies the degree-one term of the total Chern
class with that boundary \cite{SGA5}. The finite root-coordinate gluing proof in
the archive section fixes the precise form needed here. The tame
cover comparison is \emph{SGA 1}, XIII.5.2--5.6, with X.3.6 for the
discrete-valuation Abhyankar lemma. The extension of regular functions
across codimension at least two uses \emph{EGA} $\mathrm{IV}_2$,
5.8.6, 5.9.8--5.9.9 and 5.10.5, as checked below. No general FGA
representability statement supplies the missing arithmetic kernel
or the intrinsic recognition of geometric tangent bases.

The $\lambda$-ring language belongs first to $K_0$: exterior powers
define $\lambda^i$, and the determinant or first gamma-filtration
piece records line bundles. An exact locator is \emph{SGA 6}, V,
Definition 2.1, Example 2.2(b), and Lemma 2.2.1 \cite{SGA6}. Higher $K_{4k}$-groups
require the higher etale comparison used in the next section. Neither
the $K_0$ operations nor a first Chern class is that comparison map.

Crystals and divided powers concern a different coefficient setting,
especially reduction in residue characteristic $\ell$. No new theorem
about a crystal, a divided-power envelope, or a crystalline comparison
is proved in this run. The etale Kummer calculation deliberately removes
$\ell$ from the base; for noninvertible exponents one must use flat
Kummer theory and then state a separate comparison with any crystalline
or prismatic object. The inseparable counterexample in the archive
section proves why the etale assertion itself cannot be transported
unchanged. Assigning one of the present finite Selmer estimates to a
crystalline object without such a map would add an unproved assertion.

Finally, the rational mixed-Tate category over $\operatorname{Spec}
\mathbb Z$ is the category used by Deligne--Goncharov and Brown. Its
base scheme controls ramification; its rational coefficients control
which extensions remain visible. The integral arithmetic lattice of a
pro-$\ell$ image is an additional object. The counterexample at $37$
below shows that these distinctions change the answer to B3.



\clearpage
\part{Unresolved findings, failed implications, and continuation gates}
\section{Open assertions are not used as theorems}

\paragraph{OPEN-K-NUM: a specified new vanishing degree.}
The existential polynomial order bound and its explicit exponent do
not supply a small numerical constant. A new degree such as $K_{12}$
requires evaluated fixed-level cohomology bounds, or computation of the
two nonzero trace products and their gcd, or a different integral
cohomology calculation. The available primary literature checks and
the exact trace-product reduction are documented in the evidence notes.

\paragraph{OPEN-IRREG-KERNEL: irregular odd primes.}
The proved separation theorem handles cyclic even Iwasawa components of
arbitrary finite rank, including repeated and nonsplit roots. Vandiver
and odd cup-product surjectivity suffice. Published computations verify
both inputs for every odd prime below $1000$. For remaining irregular
primes, the unresolved input is these arithmetic conditions (or a
weaker noncancellation criterion); primes failing Vandiver would also
require an extension to noncyclic even modules.

\paragraph{OPEN-2-LATTICE: integral odd generation at $2$.}
The exact outer-filtration comparison and compact-group calculation
prove integral odd generation through weight $18$. The all-weight gate
is an integral Hall/Newton correction and stability theorem for the
Gamma action. No all-weight proof or counterexample is asserted.

\paragraph{OPEN-149-REL and OPEN-149-GEOM.}
The local root-coordinate obstruction is computed. The complete
relative normalizer/formal-tube comparison in the generality of the
manuscript still requires precise homotopy-exactness and descent
hypotheses. Even after that comparison, recognizing the uncompleted
global-unit image intrinsically inside all splitting classes remains
the scan-85 problem. The concrete extra local class does not prove
that such recognition is impossible.

\paragraph{Source and transcription scope.}
The broader host-field point theorem now has a closed source proof
composition, with its numerical layer bounds retained as premises;
the character theorem and rational-field specialization are also checked. Missing original
round-one B3/S9 text and missing handwritten archive images are not
reconstructed from conjecture. Mathematical editorial flags are
distinguished from verified readings in TRANSCRIPTION\_NOTES.tsv.

\section{Audit conventions and source reproducibility}
CHECK\_EVENTS.tsv records actual source readings, whole-proof checks,
independent re-proofs and exact computations. A cited theorem may be
used with its statement and hypotheses checked without a fresh audit of
its entire original proof; that distinction is retained. The source
contract accepts Alp\"oge's supplied theorems and does not reopen RH.
The pinned mirror, checksums and retrieved Lean excerpts are in the
evidence directory. Research notes also record primary URLs and exact
page or declaration locators. Earlier candidate notes are retained for
continuity but do not supersede the checked paper and latest ledger.

\begin{thebibliography}{99}

\bibitem{Kurihara1992}
M. Kurihara,
Some remarks on conjectures about cyclotomic fields and \(K\)-groups of \(\mathbb Z\),
\emph{Compositio Math.} 81 (1992), 223--236.
\url{https://www.numdam.org/item/CM_1992__81_2_223_0.pdf}.
Exact locators used: Lemma 1.2, p.225; Corollary 1.5, p.226;
Proposition 3.7, p.229; Remark 3.11, p.230;
(4), Proposition 5.1 and Remark 5.2, p.233.


\bibitem{Weibel2013}
C. A. Weibel,
\emph{The \(K\)-book: An Introduction to Algebraic \(K\)-theory},
Graduate Studies in Mathematics 145, AMS, 2013.
Author's Chapter VI:
\url{https://sites.math.rutgers.edu/~weibel/Kbook/Kbook.VI.pdf}.
Exact locators used: VI.8.2, VI.8.3, VI.9.8, VI.10.9.


\bibitem{RognesWeibel2000}
J. Rognes and C. Weibel,
Two-primary algebraic \(K\)-theory of rings of integers in number fields,
\emph{J. Amer. Math. Soc.} 13 (2000), 1--54.
\url{https://sites.math.rutgers.edu/~weibel/archive/papers-dir/RognesWeibel.pdf}.
Theorem 0.7, pp.4--5, proved as Theorem 7.11;
localization discussion p.1 and Theorem 0.6.


\bibitem{Rognes2000}
J. Rognes, \(K_4(\mathbb Z)\) is the trivial group,
\emph{Topology} 39 (2000), 267--281.
\url{https://doi.org/10.1016/S0040-9383(99)00007-5}.
Theorem 0.1, p.268.


\bibitem{DEKM2026}
M. Dutour Sikiri\'c, P. Elbaz-Vincent, A. Kupers and J. Martinet,
Voronoi complexes in higher dimensions, cohomology of
\(\operatorname{GL}_N(\mathbb Z)\) for \(N\ge8\) and the triviality of
\(K_8(\mathbb Z)\),
\emph{J. Inst. Math. Jussieu} 25 (2026), 565--590.
\url{https://doi.org/10.1017/S1474748025101394}.
Published online 24 November 2025.
Theorem 6.4; author preprint arXiv:1910.11598v1, p.19,
and Corollary 7.1, p.20:
\url{https://arxiv.org/abs/1910.11598}.


\bibitem{HartHarveyOng2017}
W. B. Hart, D. Harvey and W. Ong,
Irregular primes to two billion,
\emph{Math. Comp.} 86 (2017), 3031--3049.
\url{https://doi.org/10.1090/mcom/3211}.
Author preprint:
\url{https://arxiv.org/abs/1605.02398}.
Section 1, item (2), verifies Vandiver for every odd prime
\(\ell<2^{31}\).


\bibitem{MilneADT}
J. S. Milne, \emph{Arithmetic Duality Theorems},
second edition, BookSurge, 2006.
\url{https://www.jmilne.org/math/Books/ADT.pdf}.
Finite duality locators: I, Theorem 4.10; II, Theorem 2.9.


\bibitem{Jannsen1988}
U. Jannsen, Continuous \'etale cohomology,
\emph{Math. Ann.} 280 (1988), 207--245.
\url{https://doi.org/10.1007/BF01456052}.
Author's institutional copy:
\url{https://epub.uni-regensburg.de/26684/1/jannsen12.pdf}.
Exact inverse-limit sequence (3.1).


\bibitem{MilneANT}
J. S. Milne, \emph{Algebraic Number Theory},
course notes, version 3.08, 2020.
\url{https://www.jmilne.org/math/CourseNotes/ANT.pdf}.
Theorem 4.3, printed p.70.


\bibitem{EtingofGL2024}
P. Etingof, \emph{Lie Groups and Lie Algebras II},
MIT 18.755, Spring 2024.
Lecture 01, sections 27.2--27.3, printed pp.145--147:
\url{https://ocw.mit.edu/courses/18-755-lie-groups-and-lie-algebras-ii-spring-2024/mit18_755_s24_lec01.pdf}.
Lecture 02, section 28.1, printed pp.151--152:
\url{https://ocw.mit.edu/courses/18-755-lie-groups-and-lie-algebras-ii-spring-2024/mit18_755_s24_lec02.pdf}.
Highest-weight classification, determinant twists, and Weyl's dimension formula.


\bibitem{IshiiKummer2024}
S. Ishii,
On Kummer characters arising from Galois actions on the pro-\(p\)
fundamental groups of once-punctured CM elliptic curves,
\emph{Research in Number Theory} 11 (2025), article 38.
\url{https://doi.org/10.1007/s40993-025-00622-0}.
Author preprint arXiv:2312.04175v2:
\url{https://arxiv.org/pdf/2312.04175}.
Definition preceding Theorem 1.1 and Theorem 1.1, p.2.


\bibitem{IchimuraKaneko1989}
H. Ichimura and M. Kaneko,
On the universal power series for Jacobi sums and the Vandiver conjecture,
\emph{J. Number Theory} 31 (1989), 312--334.
Author's copy:
\url{https://www2.math.kyushu-u.ac.jp/~mkaneko/papers/04universal_power.pdf}.
Section 1.1, pp.313--314.


\bibitem{BuhlerHarvey2011}
J. P. Buhler and D. Harvey,
Irregular primes to 163 million,
\emph{Math. Comp.} 80 (2011), 2435--2444.
Author preprint:
\url{https://arxiv.org/abs/0912.2121}.
Section 4, pp.4--5 of version 2, for the Vandiver certificate criterion.


\bibitem{Kupers2019}
A. Kupers, A short proof that \(K_8(\mathbb Z)\cong0\),
author's note dated 21 October 2019.
\url{https://www.utsc.utoronto.ca/people/kupers/wp-content/uploads/sites/50/2021/01/k8zshorter.pdf}.
Theorem A and Remark 2.6, p.4.


\bibitem{BrueckEtAl2024}
B. Br\"uck, J. Miller, P. Patzt, R. J. Sroka and J. C. H. Wilson,
On the codimension-two cohomology of \(\operatorname{SL}_n(\mathbb Z)\),
\emph{Adv. Math.} 451 (2024), article 109795.
\url{https://doi.org/10.1016/j.aim.2024.109795}.
Author preprint arXiv:2204.11967v3:
\url{https://arxiv.org/abs/2204.11967}.
Theorem B: rational codimension-two vanishing for \(n\ge3\).

\bibitem{Sharifi2002}
R.~T.~Sharifi, \emph{Relationships between conjectures on the structure
of pro-$p$ Galois groups unramified outside $p$},
Proc. Sympos. Pure Math. \textbf{70} (2002), 275--284.
Author's 13-page version: \url{https://www.math.ucla.edu/~sharifi/publ/galstr.pdf}.


\bibitem{Brown2012}
F.~Brown, \emph{Mixed Tate motives over $\mathbb Z$},
Ann. of Math. \textbf{175} (2012), 949--976.
\url{https://annals.math.princeton.edu/wp-content/uploads/annals-v175-n2-p10-p.pdf}.


\bibitem{Brown2019}
F.~Brown, \emph{From the Deligne--Ihara conjecture to multiple modular values},
arXiv:1904.00179 (2019).
\url{https://arxiv.org/pdf/1904.00179}.


\bibitem{HainMatsumoto2001}
R.~Hain and M.~Matsumoto, \emph{Weighted completion of Galois groups
and Galois actions on the fundamental group of $\mathbb P^1-\{0,1,\infty\}$},
arXiv:math/0006158v2, 11 August 2001.
\url{https://arxiv.org/pdf/math/0006158}.


\bibitem{GrasJaulent2008}
G.~Gras and J.-F.~Jaulent, \emph{Note on $2$-rational fields},
J. Number Theory \textbf{129} (2009), 495--498;
arXiv:0805.1168.
\url{https://arxiv.org/pdf/0805.1168}.


\bibitem{MinacPasiniQuadrelliTan2020}
J.~Min\'a\v c, F.~Pasini, C.~Quadrelli, and Nguy{\~{\^e}}n Duy T\^an,
\emph{Koszul algebras and quadratic duals in Galois cohomology},
arXiv:1808.01695v3, 31 December 2020.
\url{https://arxiv.org/pdf/1808.01695}.


\bibitem{IshiiKernel2026}
S.~Ishii, \emph{On the kernels of the pro-$p$ outer Galois representations
associated to once-punctured CM elliptic curves},
arXiv:2312.04196v3, 6 January 2026.
\url{https://arxiv.org/pdf/2312.04196}.


\bibitem{McCallumSharifi2003}
W.~G.~McCallum and R.~T.~Sharifi,
\emph{A cup product in the Galois cohomology of number fields},
Duke Math. J. \textbf{120} (2003), 269--310.
Author's 34-page version dated 5 December 2003:
\url{https://www.math.ucla.edu/~sharifi/publ/pairing.pdf}.


\bibitem{Iwasawa1981}
K.~Iwasawa, \emph{Riemann--Hurwitz formula and $p$-adic Galois
representations for number fields},
Tohoku Math. J. \textbf{33} (1981), 263--288.
\url{https://www.jstage.jst.go.jp/article/tmj1949/33/2/33_2_263/_pdf/-char/en}.


\bibitem{SharifiIwasawaNotes}
R.~T.~Sharifi, \emph{Iwasawa Theory}, online notes, Chapter~3,
especially Corollary~3.3.14, Propositions~3.4.1 and~3.4.7,
and Corollary~3.4.8. Accessed 21 November 2026.
\url{https://www.math.ucla.edu/~sharifi/notes/iwasawa-ch03.html}.


\bibitem{DeligneGoncharov2005}
P.~Deligne and A.~B.~Goncharov,
\emph{Groupes fondamentaux motiviques de Tate mixte},
Ann. Sci. \'Ecole Norm. Sup. (4) \textbf{38} (2005), 1--56.
The supplied LaTeX source was used for \S1.6,
Propositions~1.7, 1.9 and~2.3, and \S2.4.


\bibitem{Hirose2026}
M.~Hirose, \emph{Confluence relations, Hoffman basis over $\mathbb Z_{(2)}$,
and the Deligne--Drinfeld conjecture},
arXiv:2610.23719v1, 16 October 2026.
\url{https://arxiv.org/pdf/2610.23719}.


\bibitem{Sharifi2007} R. Sharifi, \emph{Iwasawa theory and the Eisenstein ideal}, Duke Math. J. \textbf{137} (2007), 63--101. Locators refer to arXiv:math/0501236v4 (17 March 2007): Theorem 1.2 p.~2; Theorem 5.6 pp.~25--26; Theorem 5.7 and Corollaries 5.8--5.9 p.~27; weight-twelve application p.~28. \url{https://arxiv.org/abs/math/0501236}.

\bibitem{SGA1}
A.~Grothendieck et al., \emph{Rev\^etements \'etales et groupe fondamental
(SGA 1)}, Lecture Notes in Mathematics \textbf{224}, Springer, 1971;
annotated reedition, Documents Math\'ematiques \textbf{3},
Soci\'et\'e Math\'ematique de France, 2003.
Expos\'e XIII, Proposition~5.2 and Corollaire~5.3,
reedition pp.~315--316 (original marginal pp.~428--430);
relative versions, Proposition~5.5 and Corollaire~5.6;
Expos\'e X, Lemme~3.6.
\url{https://arxiv.org/pdf/math/0206203}.


\bibitem{SGA4}
M.~Artin, A.~Grothendieck, and J.-L.~Verdier (eds.),
\emph{Th\'eorie des topos et cohomologie \'etale des sch\'emas
(SGA 4)}, Tome~3, Lecture Notes in Mathematics \textbf{305},
Springer, 1973. Expos\'e IX, Th\'eorie de Kummer~3.2 and
Th\'eor\`eme~3.3; French retypeset edition, printed pp.~14--15
(PDF pages~19--20).
\url{https://www.normalesup.org/~forgogozo/SGA4/tomes/tome3.pdf}.


\bibitem{StixCuspidal2012}
J.~Stix, \emph{On cuspidal sections of algebraic fundamental groups},
author version dated 24 February 2012.
Proposition~11, p.~7; Proposition~16 and Theorem~17, p.~8;
Corollary~19, p.~9; Corollary~32, p.~12.
\url{https://www.math.uni-frankfurt.de/~stix/research/preprints/STIXcuspsec20120802homepage.pdf}.


\bibitem{EGAIV2}
A.~Grothendieck, with the collaboration of J.~Dieudonn\'e,
\emph{\'El\'ements de g\'eom\'etrie alg\'ebrique IV:
\'Etude locale des sch\'emas et des morphismes de sch\'emas,
seconde partie}, Publ. Math. IH\'ES \textbf{24} (1965), 5--231.
Th\'eor\`eme~5.8.6, p.~108; Proposition~5.9.8 and
D\'efinition~5.9.9, pp.~111--112;
Th\'eor\`eme~5.10.5 and proof, pp.~115--116.
Primary scan:
\url{https://webusers.imj-prg.fr/~leila.schneps/grothendieckcircle/EGA/EGAIV.2.pdf}.
Numdam record: \url{https://archive.numdam.org/item/PMIHES_1965__24__5_0/}.

\bibitem{DeligneWeilII}
P.~Deligne, \emph{La conjecture de Weil. II},
Publ. Math. IHES \textbf{52} (1980), 137--252.
Definitions 1.2.1--1.2.2, p.~153, and 1.2.5(iv), p.~154.
\url{https://www.numdam.org/item/PMIHES_1980__52__137_0.pdf}.


\bibitem{SGA5}
A.~Grothendieck et al., \emph{Cohomologie $\ell$-adique et fonctions $L$},
SGA~5, Lecture Notes in Mathematics 589, Springer, 1977.
Expose VII, standing invertibility convention p.~282; Proposition 3.4(i), p.~291, proof p.~292; Corollary 3.7(i), p.~295. Original scanned pages visually checked.
\url{https://wstein.org/sga/sga5/SGA5.pdf}.


\bibitem{SGA6}
P.~Berthelot, A.~Grothendieck and L.~Illusie, eds.,
\emph{Theorie des intersections et theoreme de Riemann--Roch},
SGA~6, Lecture Notes in Mathematics 225, Springer, 1971.
Expose V, Definition 2.1, Example 2.2(b), Lemma 2.2.1.
Checked section transcription:
\url{https://grothendiecksga.com/read/sga6/fr/V-2.html}.

\bibitem{KochKukkukLabute1998}
H.~Koch, S.~Kukkuk and J.~Labute, \emph{Nilpotent local class field theory},
Acta Arithmetica \textbf{83} (1998), 45--64.
Theorem~2.1, pp.~48--49 (author PDF pp.~4--5).
\url{https://www.math.mcgill.ca/labute/papers/kkl.pdf}.

\end{thebibliography}
\end{document}
